\IfFileExists{stacks-project.cls}{%
\documentclass{stacks-project}
}{%
\documentclass{amsart}
}
\usepackage{amssymb}

% For dealing with references we use the comment environment
\usepackage{verbatim}
\usepackage{fancyvrb}
\usepackage{pmboxdraw}
\newenvironment{reference}{\comment}{\endcomment}
%\newenvironment{reference}{}{}
\newenvironment{slogan}{\comment}{\endcomment}
\newenvironment{history}{\comment}{\endcomment}

% For source figures retained as assets
\usepackage{graphicx}

% For commutative diagrams we use Xy-pic
\usepackage[all]{xy}

% We use 2cell for 2-commutative diagrams.
\xyoption{2cell}
\UseAllTwocells

% We use multicol for the list of chapters between chapters
\usepackage{multicol}

% This is generally recommended for better output
\usepackage{lmodern}
\usepackage[T1]{fontenc}
\usepackage{textalpha}

% For cross-file-references
% We prefer the snapshot tmp/xref/#2.aux,
% which latexmkrc replaces atomically, so
% that documents compiled simultaneously
% never read a half-written .aux file.
% The last argument keeps links to #2.pdf.
\usepackage{xr-hyper}
\newcommand{\maybeexternaldocument}[2]{%
  \IfFileExists{#2.aux}{%
    \externaldocument[#1]{#2}%
  }{%
    \PackageWarning{stack}{Missing #2.aux; external references with prefix #1 unavailable}%
  }%
}
\let\externaldocumentorig\externaldocument
\renewcommand{\externaldocument}[2][]{%
  \IfFileExists{tmp/xref/#2.aux}{%
    \externaldocumentorig[#1]{tmp/xref/#2}[#2.pdf]%
  }{%
  \IfFileExists{#2.aux}{%
    \externaldocumentorig[#1]{#2}%
  }{%
    \PackageWarning{stack}{Missing #2.aux; external references with prefix #1 unavailable}%
  }%
  }%
}

% Package for hypertext links:
\usepackage{hyperref}

% For any local file, say "hello.tex" you want to link to please
% use \externaldocument[hello-]{hello}
\externaldocument[introduction-]{introduction}
\externaldocument[conventions-]{conventions}
\externaldocument[sets-]{sets}
\externaldocument[categories-]{categories}
\externaldocument[topology-]{topology}
\externaldocument[sheaves-]{sheaves}
\externaldocument[sites-]{sites}
\externaldocument[stacks-]{stacks}
\externaldocument[fields-]{fields}
\externaldocument[algebra-]{algebra}
\externaldocument[brauer-]{brauer}
\externaldocument[homology-]{homology}
\externaldocument[derived-]{derived}
\externaldocument[simplicial-]{simplicial}
\externaldocument[more-algebra-]{more-algebra}
\externaldocument[smoothing-]{smoothing}
\externaldocument[modules-]{modules}
\externaldocument[sites-modules-]{sites-modules}
\externaldocument[injectives-]{injectives}
\externaldocument[cohomology-]{cohomology}
\externaldocument[sites-cohomology-]{sites-cohomology}
\externaldocument[dga-]{dga}
\externaldocument[dpa-]{dpa}
\externaldocument[sdga-]{sdga}
\externaldocument[hypercovering-]{hypercovering}
\externaldocument[schemes-]{schemes}
\externaldocument[constructions-]{constructions}
\externaldocument[properties-]{properties}
\externaldocument[morphisms-]{morphisms}
\externaldocument[coherent-]{coherent}
\externaldocument[divisors-]{divisors}
\externaldocument[limits-]{limits}
\externaldocument[varieties-]{varieties}
\externaldocument[topologies-]{topologies}
\externaldocument[descent-]{descent}
\externaldocument[perfect-]{perfect}
\externaldocument[more-morphisms-]{more-morphisms}
\externaldocument[flat-]{flat}
\externaldocument[groupoids-]{groupoids}
\externaldocument[more-groupoids-]{more-groupoids}
\externaldocument[etale-]{etale}
\externaldocument[chow-]{chow}
\externaldocument[intersection-]{intersection}
\externaldocument[pic-]{pic}
\externaldocument[weil-]{weil}
\externaldocument[adequate-]{adequate}
\externaldocument[dualizing-]{dualizing}
\externaldocument[duality-]{duality}
\externaldocument[discriminant-]{discriminant}
\externaldocument[derham-]{derham}
\externaldocument[local-cohomology-]{local-cohomology}
\externaldocument[algebraization-]{algebraization}
\externaldocument[curves-]{curves}
\externaldocument[resolve-]{resolve}
\externaldocument[models-]{models}
\externaldocument[functors-]{functors}
\externaldocument[equiv-]{equiv}
\externaldocument[pione-]{pione}
\externaldocument[etale-cohomology-]{etale-cohomology}
\externaldocument[proetale-]{proetale}
\externaldocument[relative-cycles-]{relative-cycles}
\externaldocument[more-etale-]{more-etale}
\externaldocument[trace-]{trace}
\externaldocument[crystalline-]{crystalline}
\externaldocument[spaces-]{spaces}
\externaldocument[spaces-properties-]{spaces-properties}
\externaldocument[spaces-morphisms-]{spaces-morphisms}
\externaldocument[decent-spaces-]{decent-spaces}
\externaldocument[spaces-cohomology-]{spaces-cohomology}
\externaldocument[spaces-limits-]{spaces-limits}
\externaldocument[spaces-divisors-]{spaces-divisors}
\externaldocument[spaces-over-fields-]{spaces-over-fields}
\externaldocument[spaces-topologies-]{spaces-topologies}
\externaldocument[spaces-descent-]{spaces-descent}
\externaldocument[spaces-perfect-]{spaces-perfect}
\externaldocument[spaces-more-morphisms-]{spaces-more-morphisms}
\externaldocument[spaces-flat-]{spaces-flat}
\externaldocument[spaces-groupoids-]{spaces-groupoids}
\externaldocument[spaces-more-groupoids-]{spaces-more-groupoids}
\externaldocument[bootstrap-]{bootstrap}
\externaldocument[spaces-pushouts-]{spaces-pushouts}
\externaldocument[spaces-chow-]{spaces-chow}
\externaldocument[groupoids-quotients-]{groupoids-quotients}
\externaldocument[spaces-more-cohomology-]{spaces-more-cohomology}
\externaldocument[spaces-simplicial-]{spaces-simplicial}
\externaldocument[spaces-duality-]{spaces-duality}
\externaldocument[formal-spaces-]{formal-spaces}
\externaldocument[restricted-]{restricted}
\externaldocument[spaces-resolve-]{spaces-resolve}
\externaldocument[formal-defos-]{formal-defos}
\externaldocument[defos-]{defos}
\externaldocument[cotangent-]{cotangent}
\externaldocument[examples-defos-]{examples-defos}
\externaldocument[algebraic-]{algebraic}
\externaldocument[examples-stacks-]{examples-stacks}
\externaldocument[stacks-sheaves-]{stacks-sheaves}
\externaldocument[criteria-]{criteria}
\externaldocument[artin-]{artin}
\externaldocument[quot-]{quot}
\externaldocument[stacks-properties-]{stacks-properties}
\externaldocument[stacks-morphisms-]{stacks-morphisms}
\externaldocument[stacks-limits-]{stacks-limits}
\externaldocument[stacks-cohomology-]{stacks-cohomology}
\externaldocument[stacks-perfect-]{stacks-perfect}
\externaldocument[stacks-introduction-]{stacks-introduction}
\externaldocument[stacks-more-morphisms-]{stacks-more-morphisms}
\externaldocument[stacks-geometry-]{stacks-geometry}
\externaldocument[moduli-]{moduli}
\externaldocument[moduli-curves-]{moduli-curves}
\externaldocument[examples-]{examples}
\externaldocument[exercises-]{exercises}
\externaldocument[guide-]{guide}
\externaldocument[desirables-]{desirables}
\externaldocument[coding-]{coding}
\externaldocument[obsolete-]{obsolete}
\externaldocument[fdl-]{fdl}
\externaldocument[index-]{index}

% Heap Project documents
\externaldocument[k3-]{k3}
\externaldocument[torelli-]{torelli}
\externaldocument[linear-algebra-]{linear-algebra}
\externaldocument[groups-]{groups}
\externaldocument[complex-analysis-]{complex-analysis}
\externaldocument[functional-analysis-]{functional-analysis}
\externaldocument[categories-]{categories}
\externaldocument[commutative-algebra-]{commutative-algebra}
\externaldocument[ringed-spaces-]{ringed-spaces}
\externaldocument[homological-algebra-]{homological-algebra}
\externaldocument[lie-algebras-]{lie-algebras}
\externaldocument[representation-theory-]{representation-theory}
\externaldocument[smooth-manifolds-]{smooth-manifolds}
\externaldocument[differential-topology-]{differential-topology}
\externaldocument[algebraic-topology-]{algebraic-topology}
\externaldocument[differential-geometry-]{differential-geometry}
\externaldocument[algebraic-geometry-]{algebraic-geometry}
\externaldocument[symplectic-geometry-]{symplectic-geometry}
\externaldocument[complex-geometry-]{complex-geometry}
\externaldocument[hodge-theory-]{hodge-theory}
\externaldocument[%
variations-of-hodge-structures-]{%
variations-of-hodge-structures}
\externaldocument[%
homological-mirror-symmetry-]{%
homological-mirror-symmetry}
\maybeexternaldocument{deformations-}{deformations}
\externaldocument[vertex-algebras-]{vertex-algebras}
\externaldocument[springer-]{springer}
\externaldocument[infty-categories-]{infty-categories}
\externaldocument[seminars-differential-geometry-]{%
seminars-differential-geometry}
\externaldocument[seminars-complex-geometry-]{%
seminars-complex-geometry}
\externaldocument[seminars-algebraic-geometry-]{%
seminars-algebraic-geometry}
\externaldocument[seminars-analysis-pde-]{%
seminars-analysis-pde}
\externaldocument[surfaces-]{surfaces}
\externaldocument[birational-maps-commutative-algebra-]{birational-maps-commutative-algebra}
\externaldocument[cremona-transformations-]{cremona-transformations}
\externaldocument[derived-categories-of-sheaves-]{derived-categories-of-sheaves}
\externaldocument[geometric-invariant-theory-]{geometric-invariant-theory}
\externaldocument[geometric-stability-conditions-and-group-actions-]{geometric-stability-conditions-and-group-actions}
\externaldocument[moduli-spaces-of-sheaves-]{moduli-spaces-of-sheaves}
\externaldocument[bridgeland-stability-general-theory-]{bridgeland-stability-general-theory}
\externaldocument[hilbert-schemes-]{hilbert-schemes}
\externaldocument[noncommutative-abelian-surfaces-and-kummer-type-hyperkahler-varieties-]{noncommutative-abelian-surfaces-and-kummer-type-hyperkahler-varieties}
\externaldocument[mumford-tate-groups-in-hodge-theory-]{mumford-tate-groups-in-hodge-theory}
\externaldocument[hodge-birational-atoms-2026-]{hodge-birational-atoms-2026}
\externaldocument[xxii-egd-2026-]{%
xxii-egd-2026}
\externaldocument[pde-]{pde}
\externaldocument[probability-]{probability}

\externaldocument[linguistics-]{linguistics}

\externaldocument[%
16th-alga-meeting-2026-]{%
16th-alga-meeting-2026}

\externaldocument[%
iv-jornadas-lefschetz-2026-]{%
iv-jornadas-lefschetz-2026}

\externaldocument[reading-the-masters-]{%
reading-the-masters}

\externaldocument[real-analysis-]{%
real-analysis}

\externaldocument[optimal-transport-]{%
optimal-transport}

\externaldocument[history-]{history}

\externaldocument[seiberg-witten-]{seiberg-witten}

\externaldocument[history-philosophy-biology-]{history-philosophy-biology}

\externaldocument[classical-mechanics-]{classical-mechanics}
\externaldocument[electromagnetism-]{electromagnetism}
\externaldocument[waves-and-optics-]{waves-and-optics}
\externaldocument[wif-]{waves-in-fluids-iv}
\externaldocument[thermal-physics-]{thermal-physics}
\externaldocument[quantum-mechanics-]{quantum-mechanics}
\externaldocument[special-relativity-]{special-relativity}
\externaldocument[physics-]{physics}

\externaldocument[chemistry-]{chemistry}

\externaldocument[molecular-cell-biology-]{molecular-cell-biology}

\externaldocument[microbiology-]{microbiology}

\externaldocument[protistology-]{protistology}

\externaldocument[mycology-]{mycology}

\externaldocument[botany-]{botany}

\externaldocument[zoology-]{zoology}

\externaldocument[animal-histology-]{animal-histology}

\externaldocument[animal-physiology-]{animal-physiology}

\externaldocument[systematics-biogeography-]{systematics-biogeography}

\externaldocument[ecology-natural-resources-]{ecology-natural-resources}

\externaldocument[earth-sciences-]{earth-sciences}

\externaldocument[biostatistics-]{biostatistics}

\externaldocument[mathematical-systems-biology-]{mathematical-systems-biology}

\externaldocument[cimpa-el-salvador-2024-]{cimpa-el-salvador-2024}

\externaldocument[anthropology-]{anthropology}

\externaldocument[yang-mills-theory-]{yang-mills-theory}

\externaldocument[seminars-algebra-]{seminars-algebra}

\externaldocument[seminars-physics-]{seminars-physics}

\externaldocument[seminars-computing-]{seminars-computing}

\externaldocument[oktobermat-2026-]{oktobermat-2026}

\externaldocument[saet-2026-]{saet-2026}

\externaldocument[seminars-symplectic-geometry-]{seminars-symplectic-geometry}

\externaldocument[seminars-dynamics-]{seminars-dynamics}

\externaldocument[oktobermat-2024-]{oktobermat-2024}

\externaldocument[quantum-groups-]{quantum-groups}

\externaldocument[legal-2027-]{legal-2027}

\externaldocument[derived-categories-]{derived-categories}

% Theorem environments.
%
\theoremstyle{plain}
\newtheorem{theorem}[subsection]{Theorem}
\newtheorem{proposition}[subsection]{Proposition}
\newtheorem{lemma}[subsection]{Lemma}
\newtheorem{corollary}[subsection]{Corollary}

\theoremstyle{definition}
\newtheorem{definition}[subsection]{Definition}
\newtheorem{example}[subsection]{Example}
\newtheorem{exercise}[subsection]{Exercise}
\newtheorem{situation}[subsection]{Situation}

\theoremstyle{remark}
\newtheorem{remark}[subsection]{Remark}
\newtheorem{remarks}[subsection]{Remarks}

\numberwithin{equation}{subsection}

% Macros
%
\def\lim{\mathop{\mathrm{lim}}\nolimits}
\def\colim{\mathop{\mathrm{colim}}\nolimits}
\def\Spec{\mathop{\mathrm{Spec}}}
\def\Hom{\mathop{\mathrm{Hom}}\nolimits}
\def\Ext{\mathop{\mathrm{Ext}}\nolimits}
\def\SheafHom{\mathop{\mathcal{H}\!\mathit{om}}\nolimits}
\def\SheafExt{\mathop{\mathcal{E}\!\mathit{xt}}\nolimits}
\def\Sch{\mathit{Sch}}
\def\Mor{\mathop{\mathrm{Mor}}\nolimits}
\def\Ob{\mathop{\mathrm{Ob}}\nolimits}
\def\Sh{\mathop{\mathit{Sh}}\nolimits}
\def\NL{\mathop{N\!L}\nolimits}
\def\CH{\mathop{\mathrm{CH}}\nolimits}
\def\proetale{{pro\text{-}\acute{e}tale}}
\def\etale{{\acute{e}tale}}
\def\QCoh{\mathit{QCoh}}
\def\Ker{\mathop{\mathrm{Ker}}}
\def\Im{\mathop{\mathrm{Im}}}
\def\Coker{\mathop{\mathrm{Coker}}}
\def\Coim{\mathop{\mathrm{Coim}}}

% Boxtimes
%
\DeclareMathSymbol{\boxtimes}{\mathbin}{AMSa}{"02}

%
% Macros for moduli stacks/spaces
%
\def\QCohstack{\mathcal{QC}\!\mathit{oh}}
\def\Cohstack{\mathcal{C}\!\mathit{oh}}
\def\Spacesstack{\mathcal{S}\!\mathit{paces}}
\def\Quotfunctor{\mathrm{Quot}}
\def\Hilbfunctor{\mathrm{Hilb}}
\def\Curvesstack{\mathcal{C}\!\mathit{urves}}
\def\Polarizedstack{\mathcal{P}\!\mathit{olarized}}
\def\Complexesstack{\mathcal{C}\!\mathit{omplexes}}
% \Pic is the operator that assigns to X its picard group, usage \Pic(X)
% \Picardstack_{X/B} denotes the Picard stack of X over B
% \Picardfunctor_{X/B} denotes the Picard functor of X over B
\def\Pic{\mathop{\mathrm{Pic}}\nolimits}
\def\Picardstack{\mathcal{P}\!\mathit{ic}}
\def\Picardfunctor{\mathrm{Pic}}
\def\Deformationcategory{\mathcal{D}\!\mathit{ef}}


\begin{document}

\title{Seminars in Analysis and PDE}

\maketitle

\phantomsection

\label{section-phantom}

\tableofcontents

\section{Some Problems in Hamiltonian
Geometry of PDEs}
\label{section-hamiltonian-geometry-pdes}
%% speaker: Raffaele Vitolo, Università del
%% speaker: Salento
%% place: Seminário das Sextas, PUC-Rio
%% date: September 27, 2024 (repeated at IMPA
%% date: on December 5, 2024)

\medskip\noindent
{\bf Speaker.}

Raffaele Vitolo,
Universit\`a del Salento.

\medskip\noindent
{\bf Seminar information.}

Semin\'ario das Sextas, PUC-Rio,
September 27, 2024. The talk was also
presented in the Geometric Structures on
Manifolds Seminar at IMPA on
December 5, 2024.

\medskip\noindent
{\bf Abstract.}

The Hamiltonian formulation of partial
differential equations is one of the
cornerstones of the theory of integrable
systems. It is modeled on finite-dimensional
Hamiltonian systems and has an intrinsic
geometric nature. The talk reviews
differential-geometric aspects of the
Hamiltonian theory of PDEs and discusses
projective-geometric properties of
integrable systems that have emerged in
recent years.

\medskip\noindent
{\bf Hamiltonian evolutionary equations.}

Consider an evolutionary system
$$
u_t^i=f^i(t,x,u^j,u_x^j,u_{xx}^j,\ldots).
$$
A local functional has the form
$$
\mathcal{H}=\int h\,dx.
$$
Its variational derivative is
$$
\frac{\delta\mathcal{H}}{\delta u^i}
=
\sum_\sigma(-D)_\sigma
\frac{\partial h}{\partial u^i_\sigma}.
$$

\begin{definition}
\label{definition-hamiltonian-pde}
The system has a \textit{Hamiltonian
formulation} if there is a differential
operator $A$ and a local functional
$\mathcal{H}$ such that
$$
u_t^i=A^{ij}
\frac{\delta\mathcal{H}}{\delta u^j}.
$$
\end{definition}

The operator $A$ plays the role of a
Poisson tensor. It sends the differential
of a Hamiltonian functional to the
corresponding evolutionary vector field.

\begin{definition}
\label{definition-hamiltonian-operator}
A differential operator $A$ is a
\textit{Hamiltonian operator} if
$$
\{\mathcal{F},\mathcal{G}\}_A
=
\int
\frac{\delta\mathcal{F}}{\delta u^i}
A^{ij}
\frac{\delta\mathcal{G}}{\delta u^j}
\,dx
$$
defines a Poisson bracket on local
functionals.
\end{definition}

Thus $A$ must be formally skew-adjoint and
must satisfy the differential analogue of
the Jacobi identity.

There is also a purely algebraic definition
of differential operators. Let $R$ be a
ring, let $P,Q$ be $R$-modules, and let
$\Delta:P\to Q$ be additive. For $a\in R$,
write
$$
\delta_a\Delta=a\Delta-\Delta a.
$$
The operator $\Delta$ has order at most
$k$ if
$$
\delta_{a_0}\delta_{a_1}\cdots
\delta_{a_k}\Delta=0
$$
for all $a_0,\ldots,a_k\in R$. This is the
Grothendieck definition of a differential
operator.

\medskip\noindent
{\bf The KdV equation.}

\begin{example}
\label{example-kdv-bihamiltonian}
Consider the Korteweg--de Vries equation
$$
u_t=uu_x+u_{xxx}.
$$
It admits two Hamiltonian descriptions.
Set
$$
A_1=D_x
$$
and
$$
A_2=D_x^3+\frac{2}{3}uD_x
+\frac{1}{3}u_x.
$$
For the Hamiltonians
$$
\mathcal{H}_1=
\int\left(
\frac{u^3}{6}-\frac{u_x^2}{2}
\right)dx
$$
and
$$
\mathcal{H}_2=\int\frac{u^2}{2}\,dx,
$$
one has
$$
u_t=A_1
\frac{\delta\mathcal{H}_1}{\delta u}
=A_2
\frac{\delta\mathcal{H}_2}{\delta u}.
$$
\end{example}

This is the model for a bi-Hamiltonian
system.

\begin{definition}
\label{definition-compatible-hamiltonian}
Two Hamiltonian operators $A_1,A_2$ are
\textit{compatible} if every linear
combination of them is again Hamiltonian.
A system admitting such a pair is
\textit{bi-Hamiltonian}.
\end{definition}

\medskip\noindent
{\bf The Magri scheme.}

A Hamiltonian operator sends conservation
laws to symmetries. Given compatible
Hamiltonian operators, the Magri recursion
seeks functionals $\mathcal{H}_r$ satisfying
$$
A_1
\frac{\delta\mathcal{H}_{r+1}}{\delta u}
=
A_2
\frac{\delta\mathcal{H}_r}{\delta u}.
$$
Under suitable hypotheses, this produces
an infinite hierarchy of commuting flows
and conserved Hamiltonians in involution:
$$
\{\mathcal{H}_r,\mathcal{H}_s\}=0.
$$

This motivates one notion of integrability
for PDEs: the existence of infinitely many
commuting symmetries or conservation laws.
There is no direct infinite-dimensional
version of the Liouville theorem that
settles the theory, but the hierarchy may
lead to suitable coordinates or to an
inverse-scattering construction.

\medskip\noindent
{\bf Dubrovin--Novikov operators.}

A first-order homogeneous operator has the
form
$$
A^{ij}=g^{ij}(u)D_x
+b^{ij}_k(u)u_x^k.
$$
The conditions for $A$ to be Hamiltonian
have a differential-geometric meaning.
Formal skew-adjointness forces $g^{ij}$ to
be symmetric and relates $b^{ij}_k$ to a
metric connection. The Jacobi identity
then forces that connection to be
torsion-free and the metric to be flat,
when the metric is nondegenerate.

Thus the classification of these
Hamiltonian operators becomes a problem
in the geometry of flat metrics. The
classification of compatible pairs is
more rigid and encodes whole hierarchies
of conservation laws.

\medskip\noindent
{\bf Projective classification.}

The talk described a classification program
using projective reciprocal
transformations. Such transformations can
preserve the order and Hamiltonian nature
of an operator while changing its
coefficients. One first fixes a
higher-order operator and then classifies
compatible lower-order operators modulo
this group.

In the cases discussed in the talk,
second-order homogeneous Hamiltonian
operators are related to linear line
congruences, while third-order operators
and compatible first-order operators lead
to line complexes. Pl\"ucker coordinates
place these objects inside
$$
\mathbb{P}(\Lambda^2V).
$$
Compatibility of Hamiltonian operators is
thereby translated into compatibility of
projective varieties.

The broad point is that many familiar
integrable systems carry algebraic
varieties that organize their Hamiltonian
structures. Conversely, suitable
projective data can be used to construct
integrable systems. The Kaup--Broer system
was mentioned as an example arising from
this mechanism.

\section{PDE Without Tears for Geometry
Students}
\label{section-pde-without-tears}
%% speaker: Raffaele Vitolo, Università del
%% speaker: Salento
%% place: Seminário das Sextas, PUC-Rio
%% date: November 29, 2024

\medskip\noindent
{\bf Speaker.}

Raffaele Vitolo,
Universit\`a del Salento.

\medskip\noindent
{\bf Seminar information.}

Semin\'ario das Sextas, PUC-Rio,
November 29, 2024.

The geometric study of differential
equations goes back to Lie and was put
into its modern form by Ehresmann. Its
basic idea is to regard derivatives as
coordinates on a larger space.

Let $N$ have dimension $n$ and let $M$
have dimension $m$. A field is a map
$$
u:U\longrightarrow M,
\qquad U\subset N.
$$
Locally it has independent coordinates
$x^1,\ldots,x^n$ and dependent
coordinates $u^1,\ldots,u^m$. It can
equivalently be regarded as the section
$$
s:N\longrightarrow N\times M,
\qquad s(x)=(x,u(x)).
$$

\medskip\noindent
{\bf Contact and jet spaces.}

\begin{definition}
\label{definition-contact-order}
Two local sections $s_1$ and $s_2$ have
\textit{contact of order $r$} at $x$ if
their derivatives at $x$ agree through
order $r$.
\end{definition}

Thus contact of order zero means that the
sections have the same value. Contact of
order one means in addition that their
graphs have the same tangent space.

More generally, let $E$ be a manifold of
dimension $n+m$, and let $L_1,L_2\subset E$
be submanifolds of dimension $n$ through
$z$. After expressing both as graphs in a
common chart, one can compare their
derivatives. The resulting notion of
contact is independent of the chart.

\begin{definition}
\label{definition-jet-space}
The \textit{$r$-jet} of $L$ at $z$ is its
equivalence class under contact of order
$r$, denoted by $[L]^r_z$. The
\textit{jet space} is
$$
J^r(E,n)=
\bigcup_{z\in E}
\left\{[L]^r_z\right\}.
$$
\end{definition}

There are natural projections
$$
\cdots\longrightarrow J^2(E,n)
\longrightarrow J^1(E,n)
\longrightarrow E.
$$
The first jet space records tangent
$n$-planes and is therefore a
Grassmannian bundle. Higher jet spaces
record higher-order Taylor data.

Suppose now that $E\to N$ is a bundle. A
local section has coordinates
$$
(x^\lambda,u^i),
$$
and its $r$-jet has coordinates
$$
(x^\lambda,u^i,u^i_\sigma),
\qquad |\sigma|\leq r.
$$
Here $\sigma$ is a multi-index and
$u^i_\sigma$ represents the corresponding
partial derivative of $u^i$.

\begin{definition}
\label{definition-pde-jet-space}
A \textit{partial differential equation of
order $r$} is a submanifold
$$
\mathcal{E}\subset J^r(E,n).
$$
\end{definition}

A system given locally by equations
$$
F^a(x^\lambda,u^i,u^i_\sigma)=0
$$
defines such a submanifold whenever the
usual regularity conditions hold.

\begin{definition}
\label{definition-prolongation-section}
The \textit{$r$-th prolongation} of a
section $s$ is
$$
j^rs:N\longrightarrow J^r(E,n),
\qquad x\longmapsto[s(N)]^r_{s(x)}.
$$
The section is a \textit{solution} of
$\mathcal{E}$ if
$$
j^rs(N)\subset\mathcal{E}.
$$
\end{definition}

\begin{example}
\label{example-first-derivative-one}
For $n=m=1$, the equation
$$
u_x=1
$$
is a surface in $J^1(\mathbb{R}^2,1)$.
Its solutions are
$$
u(x)=x+c.
$$
Their first prolongations lie in that
surface.
\end{example}

\medskip\noindent
{\bf The inviscid Burgers equation.}

\begin{example}
\label{example-inviscid-burgers}
The Riemann, Hopf, or inviscid Burgers
equation is
$$
u_t+uu_x=0.
$$
Along a characteristic curve one has
$$
\frac{d}{dt}u(t,x(t))=0,
\qquad
\frac{dx}{dt}=u.
$$
Consequently, if the initial condition is
$$
u(0,x)=f(x),
$$
then the solution is determined implicitly
by
$$
u(t,x)=f(x-tu(t,x)).
$$

Set
$$
\xi=x-tu,
\qquad
v(t,\xi)=u(t,x).
$$
Away from points where the coordinate
change is singular, Burgers' equation
becomes
$$
v_t=0.
$$
The singularity occurs when
$$
1+tf'(\xi)=0.
$$
Thus characteristics can intersect and a
smooth solution can develop a shock in
finite time.
\end{example}

This example suggests a general strategy:
find symmetries and conservation laws, and
then seek a change of variables that
simplifies or linearizes the equation.

\medskip\noindent
{\bf The Cartan distribution.}

The images of prolonged sections determine
distinguished tangent directions on a jet
space. These directions form the geometric
structure that must be preserved by any
transformation of differential equations.

\begin{definition}
\label{definition-cartan-distribution}
The \textit{Cartan distribution} on
$J^r(E,n)$ is the distribution spanned by
tangent spaces to prolonged submanifolds,
together with the vertical directions of
the projection to $J^{r-1}(E,n)$.
\end{definition}

For a bundle $E\to N$, introduce the total
derivatives
$$
D_\lambda=
\frac{\partial}{\partial x^\lambda}
+
\sum_{i,|\sigma|<r}
u^i_{\sigma+\lambda}
\frac{\partial}{\partial u^i_\sigma}.
$$
The Cartan distribution is locally spanned
by the $D_\lambda$ and by
$$
\frac{\partial}{\partial u^i_\tau},
\qquad |\tau|=r.
$$
The total derivatives commute:
$$
[D_\lambda,D_\mu]=0.
$$
They differentiate functions on jet space
as if all jet coordinates arose from an
actual section.

\medskip\noindent
{\bf Contact transformations and
symmetries.}

\begin{definition}
\label{definition-contact-transformation}
A \textit{contact transformation} of a jet
space is a local diffeomorphism preserving
the Cartan distribution.
\end{definition}

When there is more than one dependent
variable, the Lie--B\"acklund theorem says
that every contact transformation of a
finite-order jet space is the prolongation
of a transformation of $E$. With one
dependent variable, genuine contact
transformations already occur on
$J^1(E,n)$.

Contact transformations can sometimes
linearize nonlinear equations. A model
case is the Monge--Amp\`ere equation
$$
u_{xx}u_{yy}-u_{xy}^2=1,
$$
which may be studied using a Legendre
transformation.

\begin{definition}
\label{definition-symmetry-pde}
A \textit{symmetry} of a differential
equation $\mathcal{E}$ is a contact
transformation preserving $\mathcal{E}$.
An infinitesimal symmetry is a vector field
tangent to $\mathcal{E}$ whose flow
preserves the restricted Cartan
distribution.
\end{definition}

The bracket of two infinitesimal symmetries
is again an infinitesimal symmetry. For an
ordinary differential equation, the
Lie--Bianchi theorem says that a sufficiently
large solvable Lie algebra of symmetries
allows the equation to be solved by
quadratures. For PDEs, symmetries can reduce
the number of variables or turn a PDE into
an ODE.

\medskip\noindent
{\bf Prolongations and linearization.}

Let
$$
X=\xi^\lambda
\frac{\partial}{\partial x^\lambda}
+\phi^i\frac{\partial}{\partial u^i}
$$
be a vector field on $E$. Its characteristic
is
$$
Q^i=\phi^i-u^i_\lambda\xi^\lambda.
$$
Its prolongation to the infinite jet space
can be written as
$$
X^{(\infty)}=
\xi^\lambda D_\lambda
+
\sum_{i,\sigma}D_\sigma(Q^i)
\frac{\partial}{\partial u^i_\sigma}.
$$
This separates the motion along the base
from the evolutionary, or vertical, part.

Suppose the equation is defined by
$$
F^a(x^\lambda,u^i,u^i_\sigma)=0.
$$
Its linearization is the differential
operator
$$
\ell_F(Q)^a=
\sum_{i,\sigma}
\frac{\partial F^a}
{\partial u^i_\sigma}D_\sigma(Q^i).
$$
A generalized symmetry has characteristic
$Q$ satisfying
$$
\ell_F(Q)=0
$$
on the equation and all its differential
consequences.

\begin{example}
\label{example-kdv-linearization}
For the Korteweg--de Vries equation
$$
F=u_t+uu_x+u_{xxx}=0,
$$
the linearization is
$$
\ell_F(Q)=
D_tQ+uD_xQ+u_xQ+D_x^3Q.
$$
Time translation is a symmetry, represented
by the characteristic $Q=u_t$. Verifying
this means substituting the equation and
its total derivatives into
$\ell_F(u_t)$.
\end{example}

\medskip\noindent
{\bf Conservation laws.}

\begin{definition}
\label{definition-conservation-law}
For an equation in variables $t$ and $x$,
a \textit{conservation law} is a pair of
differential functions $\rho$ and $J$ such
that
$$
D_t\rho+D_xJ=0
$$
on every solution. The function $\rho$ is
the \textit{density}, and $J$ is the
\textit{flux}.
\end{definition}

Under suitable boundary conditions,
integration gives
$$
\frac{d}{dt}\int_a^b\rho\,dx
=-J(t,b)+J(t,a).
$$
If the boundary flux vanishes, the integral
of the density is constant in time.

Geometrically, conservation laws are
horizontal differential forms that become
closed after restriction to the equation.
They are the PDE analogue of first
integrals in ordinary differential
equations and are closely tied to
symmetries and Hamiltonian structures.

\section{Regularity Phenomena in
Nonlocal Equations}
\label{section-regularity-nonlocal}
%% speaker: Makson Santos, Universidade de
%% speaker: Lisboa
%% place: Conference for the 60th anniversary
%% place: of the Department of Mathematics,
%% place: PUC-Rio
%% date: August 13, 2026

\medskip\noindent
{\bf Conference information.}

\href{%
https://sites.google.com/view/%
60thanniversary-dmat-puc-rio/Home
}{International Conference Celebrating
the 60th Anniversary of the Department
of Mathematics at PUC-Rio}

August 12--14, 2026,
RDC Auditorium, PUC-Rio,
Rio de Janeiro.

\medskip\noindent
{\bf Speaker.}

Makson Santos,
Universidade de Lisboa.

\medskip\noindent
{\bf Date and time.}

August 13, 2026, 10:30--11:30.

\medskip\noindent
{\bf Title.}

Regularity Phenomena in Nonlocal
Equations.

\medskip\noindent
{\bf Abstract.}

The talk studies regularity for solutions
of fully nonlinear nonlocal equations in
elliptic and degenerate settings. In the
elliptic case, the goal is optimal
regularity for equations whose right-hand
side lies in $L^p$, with the regularity
space determined by $p$. In the degenerate
case, one obtains at least one viscosity
solution in
$C^{1,\alpha}_{\mathrm{loc}}$ for some
$\alpha\in(0,1)$. When the order of the
operator is sufficiently close to $2$,
one also obtains H\"older estimates for
the gradient of every viscosity solution.

\medskip\noindent
{\bf Notes.} Once we have a (weak) solution
of a PDE, we study what differentiability 
properties does the function have only
because it solves the given PDE.


\section{Quantum simulation of partial
differential equations}
\label{section-liu-quantum-simulation}
%% speaker: Nana Liu
%% place: Centro Pi seminar, IMPA
%% date: November 19, 2026

\noindent
Nana Liu, Centro Pi seminar, IMPA.
Thursday, November 19, 2026, at 15:30,
Auditorium 3.

\medskip\noindent
{\bf From the announcement.}
A quantum device evolves according to
Schr\"odinger's equation. How can it
simulate a system with a different law
of evolution? Liu will describe
Schr\"odingerisation: a linear PDE or a
system of linear ordinary differential
equations is encoded in a larger system
with Schr\"odinger-type dynamics.
The construction admits implementations
with qubits, continuous quantum variables
(qumodes), or both. The announcement also
mentions viscosity solutions of nonlinear
Hamilton--Jacobi equations and applications
to Wasserstein proximal operators.
See the \href{https://impa.br/seminarios/}
{IMPA seminar listing}.

\medskip\noindent
{\bf Preparation: one temperature mode.}
Consider a thin ring parametrized by an
angle $x$. Write $u$ for the deviation of
the temperature from its mean. In suitable
units the heat equation is $u_t=u_{xx}$.
Suppose $u(x,0)=\cos x$. Since
$
\frac{d^2}{dx^2}\cos x=-\cos x,
$
a profile $u(x,t)=a(t)\cos x$ solves the
equation precisely when $a'(t)=-a(t)$.
Hence
$
u(x,t)=e^{-t}\cos x.
$
The hot spot cools and the cold spot
warms. The shape persists while its
amplitude decays. For comparison, under
the free Schr\"odinger equation
$i\psi_t=-\psi_{xx}$ the same initial
profile evolves as $e^{-it}\cos x$.
Its amplitude acquires a phase of modulus
one. Where is the decay encoded when a
larger quantum system simulates heat flow?

\medskip\noindent
{\bf Sources of these notes.}
The notes combine 53 photographs of the
slides, four audio recordings (about 88
minutes in total, transcribed with
whisper.cpp), and Daniel's own written
notes. The photographs show the room, not
the slide files, so a few lines are partly
hidden by the speaker or by the camera;
those are marked. Slides are transcribed as
shown, and the tables are copied verbatim.
Statements marked {\bf spoken} come from the
transcript. Everything marked {\bf Daniel's
note}, {\bf Daniel's question} or {\bf
Editorial answer} is not Liu's. The talk
reports results of
Jin--Liu--Yu~\cite{jly-prl} (with the
technical companion~\cite{jly-pra}) and of
Jin--Liu~\cite{jl-analog}; the slides name
these papers, and I have not compared the
formulas below with them line by line.

\medskip\noindent
{\bf Big picture.}
A quantum computer evolves a state by a
unitary, that is, by Schr\"odinger's
equation. Most differential equations one
wants to solve are not of this form: the
heat equation, for instance, makes things
decay. {\bf Schr\"odingerisation} turns any
linear PDE into a Schr\"odinger equation,
at the price of one extra dimension, the
``warped phase'' variable $\xi$. Then one
can hand the problem to quantum hardware.
For a nonlinear PDE the plan is different:
first make the problem linear, by an exact
change of variables such as the Cole--Hopf
transformation, then Schr\"odingerise. The
talk ends with the Hamilton--Jacobi
equation and Wasserstein proximal
operators, where this plan works.

\medskip\noindent
{\bf Quantum simulation.}
{\bf Slide 4.}
``Quantum simulations can efficiently
simulate quantum systems.'' An initial state
$|u(0)\rangle$ evolves to
$U(t)|u(0)\rangle=|u(t)\rangle$ where $U(t)$
is unitary, $U(t)^\dagger U(t)=1$. The slide
says that today we are reaching beyond the
limit of what classical computers can easily
simulate: quantum computers with
$O(100)$ ``quite good'' qubits, and quantum
simulators with continuous variables with
$O(1000000)$ qumodes. A system of $N$
two-level systems (qubits) corresponds to a
classical vector with $2^N$ entries.

\medskip\noindent
{\bf Slide ``Quantum simulation of ODEs and
PDEs''.}
The slide shows $D$ ODEs
$$
\frac{dX^{[k]}}{dt}=F(X^{[k]}(t)),
\qquad
X^{[k]}\in\mathbb{R}^D,
$$
with $X^{[k]}(0)=X^{[k]}_0$ and
$k=1,\dots,M$ (``molecular dynamics,
particle swarm, electrical networks\dots
etc''), and $(d+1)$-dimensional PDEs
$$
\frac{\partial u^{[k]}}{\partial t}
+F(u^{[k]},\nabla u^{[k]},
\nabla^2u^{[k]},\dots)=0,
\qquad
u^{[k]}\in\mathbb{R}^d
$$
(``fluid dynamics, financial market
behaviour, machine learning\dots etc'').
Beside them: ``To put onto quantum hardware:
How to map onto Schr\"odinger's equation in
a simple way? Schr\"odingerisation!''

\medskip\noindent
{\bf Slide ``Dynamical processes are
described by ordinary and partial
differential equations''.}
A plot of time ($1\,\mathrm{s}$, $10^{-6}$,
$10^{-10}$, $10^{-15}\,\mathrm{s}$) against
space ($1\,\text{\AA}$, $1\,\mathrm{nm}$,
$1\,\mu\mathrm{m}$, $1\,\mathrm{m}$), with
an arrow from the small and fast scale to the
large and slow one: quantum mechanics
(Schr\"odinger's equation), molecular
dynamics (Newton's equation), kinetic theory
(Boltzmann's equation), continuum theory
(Navier--Stokes equation). The box on the
slide:
$$
i\frac{d|u(t)\rangle}{dt}
=\mathbf{H}(t)|u(t)\rangle,
\qquad
|u(t)\rangle=U(t)|u(0)\rangle.
$$

\medskip\noindent
{\bf Slide ``Chief challenges''.}
``Schr\"odingerisation method: Mapping linear
PDEs to linear quantum dynamics.'' ``For
nonlinear PDEs: Formulating the nonlinear
PDE problem in terms of a linear problem,
where the dimension of the linear problem is
not too high.''

\medskip\noindent
{\bf Slide 9, a section title.}
``4 classes of ways of using quantum
hardware as simulators.''

\begin{remark}
\label{remark-liu-four-classes}
The next slide, ``Do we choose space and time
to be continuous or discrete during
simulation? Depends on the quantum
platform!'', is a two-by-two table. Its
columns are labelled ``data encoding
(Space)'': ``discrete'' and ``continuous''.
Its rows are labelled ``data processing
(Time)'': ``discrete'' (with the caption
``digital quantum computation'') and
``continuous'' (with the caption
``continuous-time quantum computation'').
The four cells read:
\begin{itemize}
\item discrete time, discrete space:
``Qubits: circuit model quantum computing;
measurement-based quantum computing'';

\item discrete time, continuous space:
``Qumodes: continuous-variable quantum
computing'';

\item continuous time, discrete space:
``Qubits: adiabatic quantum computing;
quantum walks; quantum annealing'';

\item continuous time, continuous space:
``Qumodes: quantum analogue computing''.
\end{itemize}
At the centre of the four cells the slide
writes ``quantum simulation''. The same table
reappears later, under the title
``Schr\"odingerisation is easily applicable to
not only the four main classes of quantum
devices, but also hybrid qubit-qumode
devices'', without the words ``Qubits'' and
``Qumodes''.
\end{remark}

\begin{remark}
\label{remark-liu-differential-vs-linear}
The slide ``Formulation of differential
equations versus linear algebra'' is the
table
$$
\begin{tabular}{l|l|l}
Mathematical language & Time & Space \\
\hline
Partial differential equation & Continuous
& Continuous \\
Ordinary differential equation & Continuous
& Discrete \\
Linear algebra & Discrete & Discrete
\end{tabular}
$$
The point is that the choice of hardware
(discrete or continuous in time and in space)
should match the formulation of the problem.
\end{remark}

\medskip\noindent
{\bf Daniel's note.}
``Gateway = discrete.'' How is information
embedded in the hardware? In quantum hardware
we use qubits instead of bits. Quantum is
also continuous, not only discrete:
information can be encoded continuously, and
in quantum one can use both discrete and
continuous. ``Qumode: continuous in space and
in time.'' Idea (with a question mark): find
the proper hardware to solve your problem or
perform your computation, discrete or
continuous, in time or in space?

\medskip\noindent
{\bf Editorial answer.}
I read ``Gateway'' as the gate model of
quantum computing, which is the digital row
of the table above: time is discrete, because
the evolution is chopped into gates. A
qumode is a quantum system whose state is a
wavefunction of a real variable $x$ (the
position of a particle, or the quadrature of
a light mode), so the space is continuous.
The idea is right and is the content of the
table: the slide says the choice ``depends on
the quantum platform''. Note that ``qumode:
continuous in space and in time'' is true of
the {\bf analogue} column; a digital
simulation of a qumode would still run in
discrete steps.

\medskip\noindent
{\bf Digital quantum information: the qubit.}
The slide ``Digital quantum information:
matrix language'' compares a bit,
$|0\rangle$ or $|1\rangle$, with a qubit
(for instance the ground state $|0\rangle$ and
the first excited state $|1\rangle$ of an
atomic energy level):
$$
|u\rangle=\alpha|0\rangle+\beta|1\rangle,
\qquad
\alpha,\beta\in\mathbb{C},
$$
and ``Born's rule: Probability
$|\alpha|^2+|\beta|^2=1$''. For $N$ qubits,
$|u\rangle\in\mathcal{H}^{\otimes N}$, a
$D=2^N$ dimensional Hilbert space, with
Pauli operators $\hat\sigma^{\otimes N}$.
Physical qubits: ions, spin systems,
superconducting qubits, etc.

\medskip\noindent
{\bf Digital simulation.}
The slide ``Typical digital quantum
simulation: evolution through digital
steps'': time evolution through concatenation
of a large number of gates. The aim is to
prepare
$$
|u(t)\rangle=\sum_ju_j(t)|j\rangle
=\exp(-i\mathbf{H}t)|u(0)\rangle,
$$
with $\mathbf{H}$ a Hermitian $D\times D$
matrix. Example, Trotter decomposition: if
$\mathbf{H}=\sum_l\mathbf{H}_l$, then
$$
e^{-i\mathbf{H}t}
=\prod_le^{-i\mathbf{H}_l\Delta t}
+O(\Delta t^2),
$$
and the number of steps is
$\sim t^2/\epsilon$. The picture is a
circuit (``programmable quantum simulator'',
``fault-tolerant digital quantum
simulator'').

\medskip\noindent
{\bf Analogue quantum information: the
qumode.}
The slide ``Analogue quantum information:
operator language'' (three slides in a row).
A qumode (continuous-variable, CV) is
$$
|u\rangle=\int u(x)|x\rangle\,dx,
\qquad
x\in\mathbb{R},
$$
where $u(x)$ is the wavefunction. The
position and momentum operators satisfy
$$
\hat x|x\rangle=x|x\rangle,
\qquad
\hat p|p\rangle=p|p\rangle,
$$
$$
\int|x\rangle\langle x|\,dx=1
=\int|p\rangle\langle p|\,dp,
\qquad
[\hat x,\hat p]=i\hbar\mathbf{1},
$$
``Heisenberg's uncertainty principle'', and
in green: ``Position and momentum operators
don't commute in quantum mechanics.'' Next:
$\langle p|x\rangle=\exp(-ixp)$ and the
``Fourier transform: $\mathcal{F}|x\rangle
=|p\rangle$''; ``Important correspondence we
will use later'':
$$
xu(x)\to\hat x|u\rangle,
\qquad
\frac{d}{dx}u(x)\to i\hat p|u\rangle,
$$
and in green: ``Easy to do Fourier transform
on optical systems (change of measurement
basis)!'' Finally: infinite dimensional
Hilbert space; ``Born's rule''
$\int|u(x)|^2dx=1$; and for $N$ qumodes
$$
|u\rangle=\int u(x_1,\dots,x_N)
|x_1,\dots,x_N\rangle\,dx_1\cdots dx_N,
$$
realised by ``lasers, light inside cavities,
freely moving atoms\dots etc''.

\medskip\noindent
{\bf Analogue simulation.}
The slide ``Analogue quantum simulation:
evolution in continuous time'': time
evolution through the natural dynamics of a
quantum system, by quantum control of a single
Hamiltonian in continuous time. The aim is to
prepare $|u(t)\rangle=e^{-iH_{CV}t}|u(0)\rangle$
(or $e^{-i\mathbf{H}t}|u(0)\rangle$) without
breaking up into many gates. Systems:
``quantum optomechanics, superconductors,
photonics, neutral atoms\dots etc''.

\begin{remark}
\label{remark-liu-qaq-table}
The slide titled (in red) ``$D+1$ qumode
quantum evolution solves $D$-dimensional
linear PDE!'' is a table with the columns
``Analogue (A)'' and ``Digital (D)'' and the
rows ``Quantum (Q)'' and ``Classical (C)''.
Its four cells read:
\begin{itemize}
\item Quantum, analogue: ``PDEs mapped to
continuous quantum systems. Quantum cost
linear in dimension. No discretisation
necessary'';

\item Quantum, digital: ``Discretised PDE
mapped to qubit dynamics. Quantum cost poly
in dimension and logarithmic in $N$'';

\item Classical, analogue: ``For PDEs, often
need to discretise spatially first into a
system of ODEs'';

\item Classical, digital: ``For PDEs, can have
curse-of-dimensionality''.
\end{itemize}
\end{remark}

\medskip\noindent
{\bf Daniel's note.}
It looks like the ``simulator'' is the fastest
way to compute. But what is a simulator?
{\bf Daniel's question:} what is simulation?
It looks like a cheap way to find out what is
up with your PDE without solving it.

\medskip\noindent
{\bf Editorial answer.}
A simulator is a physical system, built or
chosen so that its own law of evolution
copies the law one cares about. The trapped
atoms on the first slide are an example: an
optical lattice reproduces the Hubbard model
of electrons in a crystal, and one reads off
how the model behaves by letting it run and
measuring. A quantum computer used as a
simulator is the same idea with programmable
hardware. So a simulation does solve the
equation, but the answer is a quantum state
$|u(t)\rangle$ that encodes the whole
solution $u(t,\cdot)$ in its amplitudes, and
one cannot print it out. One can only measure
some quantities of it, with a probability of
success and a cost that depend on the
quantity (see the slides on recovering
$u$ and the observables of Hamilton--Jacobi
below). The speed-up claimed is in the number
of qubits or qumodes: $\log N$ qubits hold
$N$ grid values, and $D+1$ qumodes hold a
function of $D+1$ real variables, with no grid
at all. In that sense it is cheap, but not
because the PDE is not solved.

\medskip\noindent
{\bf Daniel's note.}
A naive question I could ask: is it true that
people design hardware to solve specific
problems?

\medskip\noindent
{\bf Editorial answer.}
Yes, and the table of the four classes is
organised by it. Analogue simulators, quantum
annealers and cold-atom lattices are built
for a given class of Hamiltonians, while
gate-based machines are programmable and
simulate by composing gates. The talk's
message is that Schr\"odingerisation does not
need to choose: the same linear problem can be
given to any of the four classes, and to
hybrids of qubits and qumodes.

\medskip\noindent
{\bf Daniel's note.}
Some basic quantum mechanics. Important:
position $\hat x$ and momentum $\hat p$, like
in classical mechanics. {\bf Must see:} they
do not commute, so they do not have
simultaneous eigenvalues; this is the
uncertainty principle, directly linked with
noncommutativity. Is it true that the Fourier
transform of position is momentum? Quantum
Fourier transform? Must recall: quantum
modes, the Born rule.

\medskip\noindent
{\bf Editorial answer.}
Yes to the first. For a wavefunction $u(x)$
the momentum wavefunction is its Fourier
transform $\tilde u(p)$, and on the basis
states this is the slide's
$\mathcal{F}|x\rangle=|p\rangle$ with
$\langle p|x\rangle=e^{-ixp}$ (with
$\hbar=1$). The commutation relation
$[\hat x,\hat p]=i\hbar$ is the reason:
$\hat p=-i\hbar\,d/dx$ in the position
picture, and then $[\hat x,\hat p]u=
i\hbar u$ by the Leibniz rule, so no state
is an eigenvector of both. The ``important
correspondence'' on the slide,
$d/dx\to i\hat p$, is this dictionary, with
$\hbar=1$. The quantum Fourier transform is the
finite-dimensional cousin of the same map: it
acts on the $2^N$ amplitudes of $N$ qubits,
by $O(N^2)$ gates. The qumode version is
what an optical system does naturally, by a
change of measurement basis, as the slide
says.

\medskip\noindent
{\bf I. Schr\"odingerisation.}
The section slide: ``Turning linear
non-Schr\"odinger's equations into
Schr\"odinger's equations''. ``Based on:
`Quantum simulation of partial differential
equations via Schr\"odingerisation',
Physical Review Letters (arXiv:2212.13969),
Shi Jin, Nana Liu*, Yue Yu; technical
companion paper Physical Review A
(arXiv:2212.14703), Shi Jin, Nana Liu*, Yue
Yu; `Analog quantum simulation of partial
differential equations', Quantum Science and
Technology (arXiv:2308.00646), Shi Jin, Nana
Liu*.'' (Parts of the last title are hidden in
the photograph.) These are \cite{jly-prl},
\cite{jly-pra} and \cite{jl-analog}.

\medskip\noindent
{\bf Warm-up with the heat equation.}
The slide ``Linear partial differential
equations: warm-up example with heat
equation'':
$$
\frac{\partial u}{\partial t}=\mathcal{L}u,
\qquad
\mathcal{L}\text{ a linear operator,}
$$
with examples ``Schr\"odinger:
$\mathcal{L}=-i(-\nabla_x^2+V(x))$'' and
``Heat with force term:
$\mathcal{L}=-(-\nabla_x^2+V(x))$''. Write
$P=-\nabla_x^2+V(x)$. For the heat equation
$u_t=-Pu$ and the operator $-P$ is
negative: the solution decays. For the
Schr\"odinger equation $u_t=-iPu$ and
$-iP$ is anti-Hermitian: the solution only
rotates. This is the whole difficulty. The
thin-ring preparation above, with
$u=e^{-t}\cos x$ against $e^{-it}\cos x$,
is this difference in the simplest case.

\medskip\noindent
{\bf The warped phase transformation.}
The slide ``Schr\"odingerisation and the
warped phase transformation'': ``Increase
problem by one dimension (`warped phase
space')''. The transformation is
$$
w(t,x,\xi)=e^{-\xi}u(t,x),
\qquad
\xi>0,
$$
the picture is the graph of $w$ against
$\xi$, and ``Recover $u(t,x)$'':
$$
u(t,x)=\int_0^\infty w(t,x,\xi)\,d\xi
\quad\text{or}\quad
u(t,x)=e^{\xi^*}w(t,x,\xi^*),
\ \xi^*>0
$$
(the second formula is partly hidden in the
photograph; it is also on the later slide
``Recovering the state $u$'').

\medskip\noindent
{\bf The equation in the warped phase space.}
The slide ``Heat equation in `warped phase
space''':
$$
\frac{\partial w}{\partial t}
=(-\nabla_x^2+V(x))\frac{\partial w}
{\partial\xi},
\qquad
w(0,x,\xi)=e^{-|\xi|}u(0,x),
$$
with $\xi\in(-\infty,\infty)$, and the text:
``Before we had $w(0,x,\xi)=e^{-\xi}u(0,x)$
for $\xi>0$. Here $w$ moves from right to
left along $\xi$, so no boundary condition
is needed at $\xi=0$. We can extend the
initial data $w(0,x,\xi)$ and $w$ to
$\xi<0$ since the solution does not impact
the region $\xi>0$. This extension to
$\xi\in(-\infty,\infty)$ is necessary for the
Fourier transform step.''

\begin{remark}
\label{remark-liu-warped-check}
Here is the one-line check that the two
slides agree. For $\xi>0$ we have
$\partial_\xi w=-w$, hence
$w_t=e^{-\xi}u_t=-Pw=P\,\partial_\xi w$. The
new equation $w_t=P\,w_\xi$ is a transport
equation in $\xi$ with speed given by the
nonnegative operator $P$, moving towards
negative $\xi$. So whatever is put on
$\xi<0$ never reaches $\xi>0$, and the
extension by $e^{-|\xi|}$ is harmless.
\end{remark}

\medskip\noindent
{\bf Fourier transform in $\xi$.}
The slide ``Take Fourier transform to make
into Schr\"odinger-like equation'':
$$
w(t,x,\xi)=\int e^{-i\xi\eta}
\tilde w(t,x,\eta)\,d\eta,
$$
and the boxed result
$$
i\frac{\partial\tilde w}{\partial t}
=\eta(-\nabla_x^2+V(x))\tilde w,
\qquad
\tilde w(0,x,\eta)=\frac{1}{\pi(1+\eta^2)}
u(0,x).
$$
``Schr\"odinger's equations for each Fourier
mode $\eta$: $t\to\eta t$.'' Indeed
$\xi\mapsto e^{-|\xi|}$ has Fourier
transform $1/(\pi(1+\eta^2))$ in this
convention, and $\partial_\xi$ becomes
multiplication by $-i\eta$, so
$\tilde w_t=-i\eta P\tilde w$. Each mode
$\eta$ is a Schr\"odinger equation with
Hamiltonian $\eta P$, which is Hermitian
because $P$ is. So the decay of $u$ has been
traded for a family of rotations, one for
each $\eta$ ($\eta$ is a new continuous
coordinate), and the rotation speeds are
different.

\medskip\noindent
{\bf Daniel's question.}
Answer in the notes and in chat: what is this
warped phase space example doing? Is she
introducing a method to solve an equation
involving quantization? I think this is the
Schr\"odingerisation itself, right?

\medskip\noindent
{\bf Editorial answer.}
Yes: the warped phase transformation is the
first step of Schr\"odingerisation, not a
separate example. The slides say so
(``Schr\"odingerisation and the warped phase
transformation''). The method has three
steps. First, add the variable $\xi$ by
$w=e^{-\xi}u$ and extend to all $\xi$. Second,
Fourier transform in $\xi$, which produces
the Schr\"odinger equation
$i\tilde w_t=\eta P\tilde w$. Third, solve that
Schr\"odinger equation on quantum hardware and
go back to $u$ (see ``recovering the state''
below). The heat equation is only the
warm-up case. There is no quantization of the
PDE in the sense of the Dirac rule: $x$ and
$\partial_x$ do become the operators $\hat x$
and $i\hat p$ (that is the
``correspondence'' slide), but this just says
how a function of $x$ is stored in a qumode.
The new idea is $\xi$.

\medskip\noindent
{\bf Continuous-variable implementation.}
The slide ``Schr\"odingerisation in
continuous-variable implementation'': with
$x\to\hat x$, $\eta\to\hat\eta$ and
$\partial/\partial x\to i\hat p$, the state
$$
|\tilde w(t)\rangle\propto\iint
\tilde w(t,x,\eta)|x\rangle|\eta\rangle
\,dx\,d\eta\propto
\exp(-iH_{CV}t)|\tilde w(0)\rangle,
$$
``Example: $H_{CV}=(\hat p^2+V(\hat x))
\otimes\hat\eta$'', and ``$\hat\eta$ can act
like any quadrature, e.g. $\hat\eta=\hat x$''.
So the $\eta$ variable is carried by one extra
qumode.

\medskip\noindent
{\bf Recovering the state $u$.}
The slide ``Recovering the state $u$'':
$u(t,x)=\int_0^\infty w(t,x,\xi)\,d\xi$ or
$u(t,x)=e^{\xi^*}w(t,x,\xi^*)$, $\xi^*>0$.
``Inverse quantum Fourier transform
$|\tilde w(t)\rangle\to|w(t)\rangle\propto
\iint w(t,x,\xi)|x\rangle|\xi\rangle d\xi$.
Recover $|u(t)\rangle\propto\int u(t,x)
|x\rangle dx$: projective measurement onto
$|w(t)\rangle$ using either
$P=\int_0^\infty|\xi\rangle\langle\xi|d\xi$ or
$P=|\xi^*>0\rangle\langle\xi^*>0|$.''
``Probability of success $\propto\|\mathbf{u}
(t)\|^2/\|\mathbf{u}(0)\|^2$. Algorithm has an
in-built way to compute normalisation
constant!'' (the last words are partly hidden.)
In words: after the evolution one measures
the extra mode and keeps the run only if
$\xi>0$. For $\xi>0$ we have
$w=e^{-\xi}u$, so on that event the state is a
product: the $x$-register holds
$|u(t)\rangle$ and the $\xi$-register holds
the fixed profile $e^{-\xi}$. The price is the
probability of keeping the run, which is the
ratio of the norms: if $u$ decays a lot (as
for the heat equation), success becomes rare.

\medskip\noindent
{\bf General methodology.}
Two slides. The first: ``Linear homogeneous
PDE; $K$th-order derivatives, $D$ spatial
dimensions'':
$$
\frac{\partial u}{\partial t}
+\sum_{k=1}^K\sum_{j=1}^Da_{k,j}
(x_1,\dots,x_D)\frac{\partial^ku}
{\partial x_j^k}+b(x_1,\dots,x_D)u=0,
$$
with ``$a_{k,j}<0$ for $k$ even'',
$\mathbf{u}(t)\equiv\int u(t,x)|x\rangle dx$,
$$
\frac{d\mathbf{u}}{dt}=-iA(\hat x_1,\dots,
\hat x_D,\hat p_1,\dots,\hat p_D)\mathbf{u},
$$
$$
A=\sum_{k=1}^K\sum_{j=1}^Da_{k,j}
(\hat x_1,\dots,\hat x_D)\,i^{k+3}\hat p_j^k
-ib(\hat x_1,\dots,\hat x_D),
$$
with the rule $x_i\to\hat x_i$,
$\partial^k/\partial x_j^k\to(i\hat p_j)^k$,
``can also have explicit time-dependence''
(partly hidden). The second slide splits
$A=A_1-iA_2$ with
$A_1=(A+A^\dagger)/2=A_1^\dagger$ and
$A_2=i(A-A^\dagger)/2=A_2^\dagger$:
``Hermitian part of $A$: Schr\"odinger-like
equation. Anti-Hermitian part of $A$:
heat-like equation. Original system stable if
$A_2>0$. We can also deal with unstable
cases too (simple modification).''
Check: $-iA=-iA_1-A_2$, so $A_1$ rotates and
$A_2$ damps, and the warped phase trick is
applied to $A_2$ only.

\medskip\noindent
{\bf Implementation scheme.}
The slide ``Implementation scheme'': ``Just
add one qumode: `Qumodisation'.'' The initial
state is $|\tilde v(0)\rangle$, made of the
ancilla $\int e^{-|\xi|}|\xi\rangle d\xi$ and
$|u(0)\rangle$; it goes through
$\exp(-i(A_2\otimes\hat\eta+A_1\otimes I)t)$;
the ancilla is projected with
$\hat P_{>0}=\int_0^\infty|\xi\rangle\langle
\xi|d\xi$, leaving
$|u(t)\rangle=\frac{1}{\|\mathbf{u}(t)\|}\int
u(t,x)|x\rangle dx$, with probability
$\|\mathbf{u}(t)\|^2/(2\|\mathbf{u}(0)\|^2)$.
``Initial ancilla state preparation'':
$|\Xi\rangle=\int e^{-|\xi|}|\xi\rangle d\xi$
is close to the Gaussian state
$|G\rangle=\frac{1}{\sqrt s\,\pi^{1/4}}
\int e^{-\xi^2/2s^2}|\xi\rangle d\xi$, with
$|\langle\Xi|G\rangle|\approx0.986$ for
$s\approx0.925$. In green: ``$D+1$ qumodes''
(cut off in the photograph).

\begin{remark}
\label{remark-liu-summary-diagram}
The summary slide, which starts from a PDE and
ends in Schr\"odingerisation, is the
following diagram. Each box is read downwards.
$$
\begin{tabular}{c}
Linear PDE \\[1ex]
$\displaystyle\frac{\partial u}{\partial t}
=\sum_kA_k\left(x,\frac{\partial^k}
{\partial x^k}\right)u$ \\[1ex]
$\Downarrow$ \\[1ex]
Embed into qumodes: \\[1ex]
$u\to\mathbf{u}\equiv\int u(t,x)|x\rangle
dx$ \\[1ex]
$\Downarrow$\quad $x\to\hat x$,\quad
$\dfrac{\partial^k}{\partial x^k}\to
(i\hat p)^k$ \\[1ex]
$\displaystyle\sum_kA_k\to-iA$,\quad
$\dfrac{d\mathbf{u}}{dt}=-iA\mathbf{u}$,
\\[1ex]
initial condition $\mathbf{u}(0)$ \\[1ex]
$\Downarrow$ \\[1ex]
Schr\"odingerisation: \\[1ex]
$\dfrac{d\tilde{\mathbf{v}}}{dt}
=-iH\tilde{\mathbf{v}}$, \\[1ex]
initial condition
$\tilde{\mathbf{v}}(0)=\int e^{-|\xi|}
|\xi\rangle d\xi\,\mathbf{u}(0)$, \\[1ex]
$H=A_2\otimes\hat\eta+A_1\otimes I=H^\dagger$
\end{tabular}
$$
In a side box: $A=A_1-iA_2$,
$A_1=\frac12(A+A^\dagger)$,
$A_2=\frac{i}{2}(A-A^\dagger)$. The bottom
line of the slide, partly cut off: ``Also
applicable to time-dependent
(non-autonomous) \dots''.
\end{remark}

\medskip\noindent
{\bf The simplest example.}
The slide ``Simplest example'': the
$D$-dimensional convection equation
$$
\frac{\partial u}{\partial t}-\sum_{j=1}^D
a_j\frac{\partial u}{\partial x_j}=0,
\quad
u(0,x)=u_0(x),
\quad
x\in\mathbb{R}^D,\ a_j\in\mathbb{R}.
$$
``Only displacement required'':
$d\mathbf{u}/dt=-iH\mathbf{u}$ with
$H=\sum_jH_j$, $H_j=a_j\hat p_j$,
$H_j=H_j^\dagger$, on $D$ qumodes. In green:
``In general, simulation of PDEs require
entangling operations and will go beyond
Gaussian operations.''

\begin{remark}
\label{remark-liu-pde-table}
The slide with the quantum simulation cost is
a table with the columns ``Equation'',
``Number of qumodes'', ``Max order term in
$H$'' and ``Number of terms in $H$''. The red
roman numerals and letters in it are section
numbers of the source paper. Its rows are
transcribed from a zoomed photograph. A red
oval enclosed everything from ``General
linear PDEs'' down to ``Maxwell''; the blue
box beside it says ``Quantum simulation cost
is only linear in dimension\dots and if use
qubits is also logarithmic in discretization
size!''
\begin{itemize}
\item {\bf General linear PDEs.}
Homogeneous (IV), 1st-order time derivative:
$D+1$ qumodes; terms
$[a_{K,j},\hat p_j^K]\otimes\hat\eta$ and
$b\otimes\hat\eta$; $O(DK)$ terms.
Inhomogeneous (IV A): add 1 more qubit;
$[a_{K,j},\hat p_j^K]\otimes\hat\eta\otimes
\sigma$ and $b\otimes\hat\eta\otimes\sigma$;
$O(DK)$. $n$-th order time derivative (IV B):
add $\log_2(n)$ more qubits; ``similar terms
to 1st-order
$\otimes\,\sigma^{\otimes\log_2n}$'';
$O(DK)$.

\item {\bf Examples.}
Liouville (V A): $D+1$;
$i[\hat p,L]\otimes\hat\eta$ and
$\{\hat p,L\}\otimes I$; $O(D)$.
Heat (V B1): $D+1$;
$\{\hat pD,\hat p\}\otimes\hat\eta$ and
$V\otimes\hat\eta$; $O(D+1)$.
Fokker--Planck (V B2): $D+1$;
$\{\hat p^2,D\}\otimes\hat\eta$,
$i[\mu,\hat p]\otimes\hat\eta$,
$\{\hat p,\mu\}\otimes I$,
$i[\hat p^2,D]\otimes I$; $O(D)$.
Black--Scholes (V B3): $2$;
$(\hat x^2\hat p^2+\hat p^2\hat x^2)\otimes
\hat\eta$; $5$.
Wave (V B4): $D+1$ + one qubit;
$\{a,\hat p^2\}\otimes\hat\eta\otimes\sigma$,
$V\otimes\eta\otimes\sigma$,
$i[a,\hat p^2]\otimes I\otimes\sigma$; $O(D)$.
Maxwell ($D=3$) (V B5): $4$ qumodes $+$ $4$
qubits; ``See Section V B5''; $30$.

\item {\bf Uncertain linear PDE} (VI).
Convection: $D+L+1$;
$i[c,\hat p]\otimes\hat\eta$ and
$\{c,\hat p\}\otimes I$; $O(DL)$.

\item {\bf Nonlinear.}
Scalar hyperbolic (VII A): $D+1$;
$i[Q,\hat\varsigma]\otimes\hat\eta$ and
$\{Q,\hat\varsigma\}\otimes I$; $O(D)$.
Hamilton--Jacobi (VII A): $2D$;
$[\hat\varsigma H,\hat p]$; $O(D)$.
$N$ nonlinear ODEs (VII B): $N+1$;
$i[Q,F]\otimes\hat\eta$ and
$\{Q,F\}\otimes I$; $O(N)$.
\end{itemize}
The glyph written $\hat\varsigma$ above looks
like a hatted final sigma (or zeta) in the
photograph and I could not confirm it.
\end{remark}

\medskip\noindent
{\bf II. Nonlinear PDEs.}
The slide ``Are nonlinear problems suitable for
quantum computation?'' has a pyramid: ``A
computation is a physical process'';
``Quantum computation is a quantum mechanical
process''; ``A quantum mechanical process is
fundamentally linear''. The next slide,
``Making nonlinear problem look linear'',
has two arrows: ``Approximations'' and
``Exact mapping''. The exact mapping used in
the rest of the talk is the Cole--Hopf
transformation.

\begin{remark}
\label{remark-liu-nonlinear-pdes}
The slide ``Nonlinear PDEs: Strongly nonlinear
and weakly dissipative regime'', verbatim
(it has three pictures: a vortex street, a
cloud billow and a traffic jam). The first
letters of the last two items are hidden by the
speaker in the photograph and are restored
from the words that are visible.
\begin{itemize}
\item {\bf Nonsmooth optimal control:}
Competing optimal trajectories, switching
surfaces, front collisions, reachability, and
differential games

\item {\bf Far-from-equilibrium transport:}
Strong aggregation, multimodal densities,
metastable transitions, and nonlinear
data-assimilation updates

\item {\bf Shocks and caustics:}
Ray crossing, wave breaking, traffic jams,
hydraulic jumps, and compressible-flow shocks

\item {\bf Turbulence and microstructure:}
Intermittent fluid cascades, unresolved
oscillations, phase mixtures, and
sharp-interface limits

\item {\bf Sharp-front diffusion:}
Porous-medium invasion, combustion fronts,
thermal runaway, and persistent localized hot
spots

\item {\bf Coherent structures and defects:}
Solitons, rogue waves, vortices, domain
[the end of this line is hidden] and
topological-defect dynamics
\end{itemize}
\end{remark}

\begin{remark}
\label{remark-liu-solution-types}
The slide ``Different solution types:
different applications and different
methods'', verbatim.
\begin{itemize}
\item {\bf Weak (distributional) solution:}
presence of jump discontinuities (e.g. shock
waves in compressible gas)

\item {\bf Viscosity solution:} remains
single-valued and continuous but develops
corners at shock surfaces\dots etc, but
selects unique non-smooth solution (e.g.
optimal control)

\item {\bf Multi-valued solution:} several
physically distinct solutions (e.g.
geometric optics and seismology). {\it Quantum
algorithms on multivalued solutions:
arXiv:2202.07834}

\item {\bf Measure-valued solution:} turbulent
or highly oscillatory limits of nonlinear
fluid equations. {\it Quantum algorithm
papers on Young-measure solutions:
arXiv:2604.11825}
\end{itemize}
``Different types of solutions encode
different physically or analytically
important admissibility conditions.''
(I have not checked the two arXiv numbers; the
second one is a very recent preprint.)
\end{remark}

\medskip\noindent
{\bf Daniel's note.}
Epiphany: distributions (from the Kodaira
lectures) are weak solutions! Not functions:
they make sense only under the integral sign.
This is gold.

\medskip\noindent
{\bf Editorial answer.}
The idea is right, and Liu says the same in
the talk (see the spoken remarks on
solutions), with one refinement.
A distribution is defined by how it acts on
test functions, as in Definition
\ref{%
hodge-theory-definition-distribution-board}
of the Hodge theory notebook, and
derivatives of distributions are defined by
moving them onto the test function. A
distributional (or weak) solution is a
function or distribution that satisfies the
equation after pairing with test functions,
so the equation need not make sense point by
point. But a weak solution is not
necessarily not a function: every locally
integrable function is a distribution, and
the shock waves on the first line of the
slide are ordinary piecewise smooth
functions with jumps, which solve the
equation only in this weak sense. A Dirac
delta is the other kind. The viscosity
solution of the next line is a different
notion: it is a continuous function, and the
equation is tested with smooth functions that
touch it from above and from below, instead of
being integrated by parts.

\medskip\noindent
{\bf Daniel's question.}
Should I email her about something? She said
to email her about the other types of
solutions.

\medskip\noindent
{\bf Editorial answer.}
Her remark was an invitation (she says, ``you
can come and ask me, or just email me''), so a
short email is natural. A good first message would ask
for the references on quantum algorithms for
multivalued and Young-measure solutions (the
two arXiv numbers on the slide) and say which
equation you care about. I will not send
anything.

\medskip\noindent
{\bf III. The Hamilton--Jacobi equation.}
The slide ``Hamilton--Jacobi equation'':
$$
\frac{\partial S}{\partial t}
+\mathcal{H}(t,x,\nabla S)=0,
\qquad
S_0(x)=S(0,x),
$$
``$\mathcal{H}$ is continuous''. ``Applications
in classical mechanics, geometric optics,
mean-field games, optimal control, machine
learning\dots etc.''

\medskip\noindent
{\bf Viscosity solutions.}
The slide ``Hamilton--Jacobi PDE: Viscosity
solutions'': ``{\bf Intuition:} when the
derivative of the solution does not exist,
replace it by a (unique) smooth function that
`touches above and below' the solution.''
``Viscosity sub-solution: For every smooth test
function $\phi\in C^1$, if $S-\phi$ has a
local maximum at $(t_0,x_0)$, then
$$
\frac{\partial S(t_0,x_0)}{\partial t}
+\mathcal{H}(t_0,x_0,\nabla\phi(t_0,x_0))
\le0.
$$
Viscosity super-solution: For every smooth test
function $\phi\in C^1$, if $S-\phi$ has a
local minimum at $(t_0,x_0)$, then
$$
\frac{\partial S(t_0,x_0)}{\partial t}
+\mathcal{H}(t_0,x_0,\nabla\phi(t_0,x_0))
\ge0.\text{''}
$$
``A viscosity solution is defined as a
solution that satisfies both the viscosity
sub-solution [and super-solution] and
satisfies the given initial/boundary data''
(a few words are hidden by the speaker).

\medskip\noindent
{\bf Daniel's question.}
Tell me in chat: what is viscosity?

\medskip\noindent
{\bf Editorial answer.}
In a fluid, viscosity is the internal friction
that smooths out velocity differences, for
instance in honey. The model equation is
Burgers' equation
$u_t+uu_x=\nu u_{xx}$, with $\nu>0$ the
viscosity: it has smooth solutions, while for
$\nu=0$ the solutions develop shocks in finite
time. A {\bf viscosity solution} of
a Hamilton--Jacobi equation is the limit, as
$\nu\to0$, of the solutions of the equation
with the term $\nu\Delta S$ added. The
definition on the slide recognises that limit
without computing it: at a point where $S$
has no derivative, one tests with smooth
functions touching $S$ from above (then $S$
is a sub-solution there) and from below
(super-solution). The point of the
adjective is uniqueness: among all the
non-smooth functions that solve the equation
almost everywhere there are many, and under
the usual assumptions exactly one is the
viscosity solution.

\medskip\noindent
{\bf Viscous Hamilton--Jacobi and
Cole--Hopf.}
The slide ``Viscous Hamilton--Jacobi PDE,
Quadratic Hamiltonian'' (shown twice):
``We begin with a special example of the
quadratically nonlinear viscous
Hamilton--Jacobi PDE
$$
\frac{\partial S_\nu}{\partial t}
+\mathcal{H}(t,x,\nabla S_\nu)
=\nu\sum_{j=1}^d\frac{\partial^2S_\nu}
{\partial x_j^2},
$$
$$
\mathcal{H}(t,x,\nabla S_\nu)
=\sum_{j=1}^d\frac12\left(\frac{\partial
S_\nu}{\partial x_j}\right)^2+V(t,x),
\qquad
S_\nu(0,x)=S_0(x).\text{''}
$$
``Define a Cole--Hopf transformation with
normalization constant:
$$
S_\nu(t,x)=-2\nu\ln(\mathcal{N}_0u(t,x)),
\qquad
\mathcal{N}_0^2=\int e^{-S_\nu(0,x)/\nu}dx.
\text{''}
$$
Next slide: ``Then it is simple to verify
that $u(t,x)$ satisfies the $d$-dimensional
heat equation with a linear source term:
$$
\frac{\partial u}{\partial t}
-\frac{V(t,x)}{2\nu}u
=\nu\sum_{j=1}^d\frac{\partial^2u}
{\partial x_j^2},
\quad
u_0(x)=u(0,x),
\quad
\int|u_0(x)|^2dx=1.
$$
Defining $R_\nu=\nabla S_\nu$, we get
multidimensional Burgers' equation (for
curl-free solution):
$$
\frac{\partial R_\nu}{\partial t}
+R_\nu\cdot\nabla R_\nu+\nabla V
=\nu\sum_{j=1}^d\frac{\partial^2R_\nu}
{\partial x_j^2}.\text{''}
$$
Note that the exponent in the slide's formula
$u=e^{-S_\nu/(2\nu)}/\mathcal{N}_0$ (the
inverse of the logarithm) gives
$u_t=\nu\Delta u+\frac{V}{2\nu}u$, which is
exactly the slide's heat equation. So the
quadratic viscous Hamilton--Jacobi equation is
linear in disguise, and the linear problem is
what Schr\"odingerisation handles (the
source term $V/2\nu$ plays the role of the
$b$ of the general methodology).

\medskip\noindent
{\bf Errors at finite viscosity.}
The slide: ``If $S(t,x)$ is the viscosity
solution of the HJ PDE and $S_\nu(t,x)$ is the
solution of the viscous HJ PDE:
$$
\sup_{0\le t\le T}\sup_{x\in\mathbb{T}^d}
|S(t,x)-S_\nu(t,x)|\le C\nu^\beta,
$$
$$
\sup_{0\le t\le T}\|\nabla S(t,x)
-\nabla S_\nu(t,x)\|_{l_2(\mathbb{T}^d)}
\le C\nu^{\beta/2}.\text{''}
$$

\medskip\noindent
{\bf Beyond quadratic Hamiltonians.}
Two slides, ``Alternative formulation
needed''. ``{\bf Intuition of what we want:}
a formulation of HJ in terms of taking small
time-steps (where in each small step the
Cole--Hopf works).''
$$
S_D^{n+1}(x)=G[S_D^n](x)
=\inf_\gamma\int\bigl(h\mathcal{L}(x,v)
+S_D^n(x+hv)+2\nu\ln\gamma(v)\bigr)
\gamma(v)\,dv,
$$
``where an entropy regulariser
$\gamma(v)\ln\gamma(v)$ is added for
$\nu\ne0$; $\gamma(v)$ is a probability
density over $\mathbb{R}^d$ so
$\int_{\mathbb{R}^d}\gamma(v)dv=1$.'' ``Why
does this work? Rough idea (ignoring
entropy regularization we get non-viscous
HJ):
$$
\partial S_D^n/\partial t\sim
(S_D^{n+1}-S_D^n)/h=(G[S_D^n]-S_D^n)/h
$$
$$
\sim\frac{h\mathcal{L}(x,v)+S_D^n(x+hv)
-S_D^n(x)}{h}
\sim\inf_v\mathcal{L}(x,v)+h\nabla S_D\cdot v
$$
$$
\sim\inf_v\mathcal{L}(x,v)-p\cdot v
\sim-\mathcal{H}(x,p)
\quad\Longrightarrow\quad
\frac{\partial S}{\partial t}
+\mathcal{H}=0.\text{''}
$$
(The photographs of these two slides carry
the label ``Page 50 of 118'' of the PDF.)
The second slide: ``Entropy regularization
helps in two ways: It is the correct term in
the presence of viscosity and need to add
uncertainty, where viscosity is acting like
inverse temperature. With this entropy
regularizer, can find the exact solution of the
minimizer: solution to a `free-energy
minimization' problem is a Gibbs state.''

\medskip\noindent
{\bf Daniel's question.}
Let's have a look at Hamilton--Jacobi. The
speaker has lots of work on this. Is it the
Lagrangian formulation of mechanics?

\medskip\noindent
{\bf Editorial answer.}
Not quite, but they are close relatives. The
Lagrangian formulation of mechanics describes
motion by minimising the action
$\int\mathcal{L}(x,\dot x)\,dt$, whose
critical paths solve the Euler--Lagrange
equations. The Hamilton--Jacobi equation
describes the {\bf minimum value} of that
action, viewed as a function $S(t,x)$ of
the endpoint (the value function), and its
gradient $p=\nabla S$ is the momentum. The
Hamiltonian $\mathcal{H}(x,p)$ is the
Legendre transform of the Lagrangian,
$\mathcal{H}(x,p)=\sup_v(p\cdot v-\mathcal{L}
(x,v))$, which is exactly the last step of
the slide's heuristic ($\inf_v\mathcal{L}-p
\cdot v=-\mathcal{H}$). The slide's
time-stepping formula $S^{n+1}=\inf\int
(h\mathcal{L}+S^n(x+hv)\dots)$ is the
dynamic programming principle of this
minimisation problem. In symplectic
terms, the graph of $\nabla S$ is a
Lagrangian submanifold of the cotangent
bundle, which is why classical solutions
break down when caustics appear.

\medskip\noindent
{\bf Recovering the observables of HJ.}
The slide ``Recovering observables of HJ,
Using amplitude-encoded quantum state of
linear PDE'': ``Once we have $|u(t)\rangle$,
then what? Some examples of quantities of
interest: (1) the value of $S$ at a point,
(2) the gradient $\nabla S$ at a point, (3) the
minimum value of $S$, (4) the value of a
known function $f(x)$ at the location of the
minimum of $S$. We will present here the
continuous-variable version of the protocols
for simplicity, but digital versions are also
available'' (the photograph hides a few
letters). The photographs of this slide carry
the label ``Page 54 of 118''.

\medskip\noindent
{\bf Value of the HJ solution.}
The slide ``Value of HJ solution'': ``We first
observe that at a location $x=x_a$, if we
measure the probability
$|\langle x_a|u(t)\rangle|^2$, then we can
write
$$
-\nu\ln|\langle x_a|\mathbf{u}(t)\rangle|^2
=-\nu\ln\left(\frac{e^{-S_\nu(t,x_a)/\nu}}
{\|\mathbf{u}(t)\|^2}\right)
=S_\nu(t,x_a)+\nu\ln\|\mathbf{u}(t)\|^2.
\text{''}
$$
``Schr\"odingerisation can be used to directly
measure the normalization constant for
$|u(t)\rangle$. Cost order
$O(d^4/\epsilon^2)$, but can be improved.''

\medskip\noindent
{\bf Minimum value of the HJ solution.}
The slide ``Minimum value of HJ solution'':
``{\bf Assumption.} For the function
$S_\nu(t,x)$ at some constant $t$, we will
assume the following properties: (1)
$x_\nu^*$ is the unique global minimiser of
$S_\nu(t,x)$; (2) $x_\nu^*$ is an interior
point in the domain; (3) The Hessian
$H_{S_\nu}(t,x_\nu^*)=\nabla^2S_\nu(t,x_\nu^*)$
is positive definite (so we have a minimum at
that point); (4)
$\|\mathbf{u}(t)\|_{CV}^2=\int e^{-S_\nu(t,x)
/\nu}dx<\infty$ since we assume $|u(t)
\rangle_{CV}$ is a physical state. No extra
assumption is needed in the discrete-variable
case since the normalisation constant is
always bounded.'' ``In the small
$\nu\le O(1/d)\ll1$ regime, it is possible to
show $|S_{\nu,\min}+\nu\ln\|\mathbf{u}(t)
\|^2|\lesssim\tilde O(\nu d)$.''
``Schr\"odingerisation can be used to directly measure the
normalization constant for $|u(t)\rangle$
[algorit]hm and cost as before.'' The remaining
observables (2) and (4) are not on the
photographs I have.

\medskip\noindent
{\bf IV. Wasserstein proximal operators.}
The slide ``Wasserstein proximal operators,
Applications'': ``{\bf Question:} which new
distribution best reduces the objective while
minimizing transport away from the current
distribution? Generative learning (deep
connection between Wasserstein proximal
operators and diffusion models); Mean-field
games (Wasserstein proximal operator as
fundamental building block for computing
equilibrium dynamics of large interacting
systems); Stochastic optimal control
(connection between Wasserstein proximal
operator and Schr\"odinger bridge problem).''

\medskip\noindent
{\bf The proximal operator for probability
distributions.}
``Proximal step for vectors (e.g. gradient
descent) for minimizing a function $F(x)$:
$$
x_{k+1}=\arg\min_x\left\{\frac{1}{2\tau}
\|x-x_k\|^2+F(x)\right\}.
$$
Suppose one wants to minimize instead a
functional $F(\rho)$ of a probability
distribution $\rho$, then need a proximal step
that actually preserves probability
normalisation and positivity:
$$
\rho_{k+1}=\arg\min_\rho\left\{\frac{1}{2\tau}
\mathcal{W}^2(\rho,\rho_k)
+\mathcal{F}(\rho)\right\}.\text{''}
$$
Then ``Wasserstein distance'': ``Want a
definition of penalty $W^2(\rho,\rho_k)$ to
denote the cost of transforming one
probability distribution $\rho$ into another
probability distribution $\rho_k$. A dynamical
formulation of Wasserstein-2 distance:
$$
W_2^2(p_0,p_1)=\min_{u(t,z),\rho(t,z)}
\int_0^1\!\!\int\rho(t,z)\|u(t,z)\|^2
dz\,dt,
$$
subject to
$$
\frac{\partial\rho(t,z)}{\partial t}
+\nabla\cdot(\rho(t,z)u(t,z))=0,
\ \ \rho(t=0,z)=p_0(z),
\ \ \rho(t=1,z)=p_1(z).\text{''}
$$
Then the ``Wasserstein proximal operator
(WPO)'':
$$
\rho_T:=\mathrm{WProx}_{T,\mathcal{F}}
(\rho_0):=\arg\min_{q\in\mathcal{P}_2
(\mathbb{R}^d)}\mathcal{F}(q)
+\frac{\mathcal{W}(\rho_0,q)^2}{2T}.
$$
``Suppose we focus on the linear energy
functional $\mathcal{F}(q)=\int_{\mathbb{R}^d}
V(x)q(x)dx$. We include a small diffusion term
in the transport term (regularized Wasserstein
proximal operator):
$$
\rho_T=\mathrm{WProx}_{T,\mathcal{F},\beta}
(\rho_0):=\arg\min_{\rho,v,q}
\int_0^T\!\!\int_{\mathbb{R}^d}\frac12
\|v(t,x)\|^2\rho(t,x)\,dx\,dt
+\int_{\mathbb{R}^d}V(x)q(x)\,dx
$$
subject to the Fokker--Planck equation
$$
\partial_t\rho(t,x)+\nabla\cdot(\rho(t,x)
v(t,x))=\beta\Delta\rho(t,x),
\ \ \rho(0,x)=\rho_0(x),
\ \ \rho(T,x)=q(x),
$$
where $\beta>0$ is the small regularization
parameter, the viscosity coefficient.''

\medskip\noindent
{\bf Solving for the optimization problem.}
The slide ``Lagrange multiplier
formulation'': ``Optimality conditions and
the forward-backward system. To derive the
optimality conditions, we introduce a
Lagrange multiplier $\Phi(t,x)$ to enforce
the Fokker--Planck constraint. The Lagrangian
is
$$
\mathcal{L}(\rho,v,\Phi)=\int_0^T\!\!
\int_{\mathbb{R}^d}\Bigl(\tfrac12\|v\|^2\rho
+\Phi\bigl(\partial_t\rho+\nabla\cdot(\rho v)
-\beta\Delta\rho\bigr)\Bigr)dx\,dt
+\int_{\mathbb{R}^d}V(x)\rho(T,x)\,dx.
$$
Taking variations with respect to $\rho$ and
$v$ yields the optimality conditions. In
particular, variation with respect to $v$
gives $v=\nabla\Phi$.'' The next slide,
``Euler--Lagrange equations'': ``Euler--Lagrange
equations gives the forward-backward PDE
system (Fokker--Planck and Hamilton--Jacobi
PDEs)''
$$
\partial_t\rho+\nabla\cdot(\rho\nabla\Phi)
=\beta\Delta\rho,
\qquad
\partial_t\Phi+\tfrac12|\nabla\Phi|^2
=-\beta\Delta\Phi,
$$
$$
\rho(0,x)=\rho_0(x),
\qquad
\Phi(T,x)=-V(x),
$$
``Proximal operator is then the solution
$\rho_T:=\mathrm{WProx}_{T,\mathcal{F},\beta}
(\rho_0)=\rho(T,x)$.''

\medskip\noindent
{\bf Daniel's question.}
Another thing I should know and do not know:
what is a Lagrange multiplier?

\medskip\noindent
{\bf Editorial answer.}
It is the standard device for minimising under
a constraint. To minimise $f(z)$ subject to
$g(z)=0$ one looks for critical points of
$f+\lambda g$ in all variables, including the
new variable $\lambda$ (the multiplier);
varying $\lambda$ returns the constraint, and
varying $z$ gives $\nabla f=-\lambda\nabla g$:
at the optimum the gradient of $f$ is
parallel to the gradient of the constraint.
Here the unknown is a pair $(\rho,v)$, the
constraint is the Fokker--Planck equation at
every point $(t,x)$, so there is one multiplier
for each point: a function $\Phi(t,x)$. The
variation in $v$ of $\frac12\|v\|^2\rho
+\Phi\nabla\cdot(\rho v)$ gives, after
integrating by parts, $\rho v-\rho\nabla\Phi=0$,
that is the slide's $v=\nabla\Phi$; the
variation in $\rho$ then gives the
Hamilton--Jacobi equation for $\Phi$ with
terminal value $-V$. So the multiplier
is not a bookkeeping device only: it becomes
the velocity potential of the optimal
transport, and the Hamilton--Jacobi equation
appears here as the equation of the
multiplier.

\medskip\noindent
{\bf Solving the forward-backward system.}
Two slides. ``Making everything linear first'':
``In Fokker--Planck equation, need to remove
Lagrange multiplier $\Phi$ out of dynamics
into initial condition (otherwise hard to
simulate in quantum way)'', the same system as
above, and ``Use Cole--Hopf again!''
$$
\psi(t,x)=\rho(t,x)e^{-\Phi(t,x)/(2\beta)}
\quad\Longrightarrow\quad
\partial_t\psi=\beta\Delta\psi,
$$
$$
\psi(0,x)=\rho_0(x)e^{-\Phi(0,x)/(2\beta)}
=\rho_0(x)e^{S_\beta(T,x)/(2\beta)}
=\frac{\rho_0(x)}{\eta(T,x)},
$$
(the denominator is partly hidden) and the
final solution
$\rho(T,x)=\psi(T,x)e^{-V(x)/(2\beta)}
=\psi(T,x)\eta_0(x)$. The last slide I have,
``Solving forward-backward PDE system:
Quantum simulation'', repeats the system
and the final solution: ``Original system to
solve for final $\rho(T,x)$ which is the
proximal operator''. The photographs stop
there.

\medskip\noindent
{\bf Spoken: the talk itself.}
The long recording (72 minutes) is a rough
automatic transcript, with errors in names and
formulas; I use it only for what Liu says in
addition to the slides, in the order of the
talk. The host introduces her and she thanks
him; it is her first visit to South America
and to Rio. She says the audience shares an
interest in PDEs, and she asks them to
interrupt.
\begin{enumerate}
\item {\bf What quantum simulation means.}
It means letting a quantum system evolve, in
the closed case by unitary dynamics, which are
reversible; reversibility will matter. On a
laptop one has to transform the quantum
dynamics first; here the device itself works
with quantum states. She says that today we
are beyond what classical computers easily
simulate, for example a thousand strongly
interacting atoms, and that a quantum system
with a thousand atoms can simulate that
dynamics. Qubits: hundreds today; continuous
variable systems (light): millions of modes.
Linear quantum dynamics is Schr\"odinger's, so
the question is how to put other equations
into that form.

\item {\bf The four classes.} In classical
computing the data processing is in
discrete steps and the encoding is by bits
(the top left corner of the table), and a
talk on quantum computing usually also stays
there (circuit model, measurement-based
model). But there were analogue classical
devices before general computers, such as the
slide rule, which encodes numbers
continuously, and continuous time is the more
natural notion for physics: discrete time is
a special case of it. Her point about the
table of the formulations (PDE, ODE, linear
algebra): with hardware that is not yet good,
the physical device must match the
mathematical language as closely as possible.
A PDE has continuous time and continuous
space, so one wants a platform of that kind.

\item {\bf Qubits.} The phase matters:
a superposition is complex, not a classical
probability. For the Trotter step she mentions
other decompositions. The difficulty is that
errors propagate over long evolutions, so one
needs fault-tolerant devices, whose cost of
correcting errors may exceed the classical
simulation.

\item {\bf Qumodes.} A rule she repeats: for
quantum, whatever you have in classical
mechanics, put hats on. The Fourier transform
between $x$ and $p$ is the basis of the
quantum Fourier transform used in Shor's
algorithm and others. For the correspondence
$d/dx\to i\hat p$ she calls it ``a little
exercise'' and notes that $x$ and $d/dx$ do not
commute either, exactly as $\hat x$ and
$\hat p$ do, so everything is consistent. A
single mode is, for instance, the light of a
laser pointer (a coherent state), and $N$ modes
are $N$ beams, though they need not be
photonic.

\item {\bf The table of costs.} In the
classical corners one first discretises, and
the cost has the curse of dimensionality. The
quantum digital corner: polynomial in the
dimension and logarithmic in the
discretisation. The quantum analogue corner:
linear in the dimension, no discretisation.
She says this needs entanglement.

\item {\bf The warm-up.} The heat and
Schr\"odinger equations differ by one factor
$i$. One may absorb it in ``imaginary time'',
but then one cannot map back to real time.
Instead, as the Argand plane makes the
imaginary numbers real by adding a dimension,
one adds a real dimension, a kind of
dilation. Heat diffuses and loses
information, while Schr\"odinger evolution is
reversible, so the target should be an
equation where one can go forward and
backward. The name ``warped'' comes from her
past as a high-energy physicist: the idea of
a warped extra dimension in which gravity
leaks. After the transformation, with the
operator treated as a constant, the equation
in $\xi$ is a translation, which does not
change shape; then the Fourier transform turns
the derivative into the missing factor $i$,
and ``$\eta$ is a real number'', so each
Fourier mode is a Schr\"odinger equation with
dilated time. Then she remarks that one does
not need to discretise again: in the
continuous-variable picture $\eta$ is just
another mode, and the Hamiltonian acts on two
modes through a tensor product, an entangling
operation, which classical systems do not
have.

\item {\bf Recovery.} The inverse Fourier
transform is a change of basis (from the
$\eta$ basis to the $\xi$ basis), and the
integral over $\xi>0$ is a measurement of the
extra mode in which one keeps only $\xi>0$; if
it fails, ``we try again''. The success
probability is related to the normalisation.

\item {\bf General method.} Hermitian part:
already Schr\"odinger-like, nothing to do.
The heat equation is the case where it is
zero. The extra qumode only needs the
ancilla state, close to a Gaussian and to a
coherent state, which has already been made in
the lab, and it is the same for every PDE.
She calls the method a meta-algorithm which
applies to any of the four corners, to the
choice of what to discretise, and to hybrid
qubit-qumode devices. One goes back from a
Hamiltonian to the PDE by reading off $A_1$,
$A_2$, then $A$. Experiments: the convection
equation needs only a displacement operation,
which is easy in quantum optics, and was done
with collaborators (the transcript says up to
``20,000 words'', that is, probably
dimensions; I could not verify the number or
the collaborator's name). For a Fokker--Planck
equation with linear drift and additive
noise the term $\hat p^2\otimes\hat\eta$ is
non-Gaussian and not easily available in
optics, so her group works with a
superconducting team on chips; several have
been built and redesigned, so this is ongoing
experimental work. If one finds a system with
the Hamiltonian naturally, the simulation is
one gate.

\item {\bf Nonlinear PDEs.} Quantum mechanics
is linear, so one must turn the problem
linear, and there are two approaches.
Linearisation (Carleman type, or closure
methods) gives an infinite linear system that
one truncates; then either the dimension is
still too high (``it is not magical'') or the
truncation works only for short times, weak
nonlinearity or strong dissipation, where
nonlinear effects die out, which is not where
the interesting behaviour is. For strongly
nonlinear, weakly dissipative regimes one
wants an exact mapping, which in general may
not exist and is open. So her group studies
special PDEs. She names reinforcement
learning, optimal control, transport, shocks
and caustics and turbulence as the examples.
\end{enumerate}

\medskip\noindent
{\bf Spoken: on solutions.}
What does ``solution'' mean for a nonlinear PDE?
The classical notion is appropriate only for
short times, without shocks or caustics. The
weak notion: ``a famous example is the Dirac
delta function; it is not really a function, it
is a distribution, because it only makes sense
under an integral sign. Similarly the
solution of the PDE only makes sense under an
integral sign.'' (My paraphrase of the
transcript.) A viscosity solution is a
continuous function and is characterised by
smooth test functions touching it from above
(and below); it is the right notion for
optimal control. There are also multivalued
solutions and measure-valued ones (one asks for
the probability that the solution takes a
value). Different types of solutions encode
different conditions and need different
algorithms. She will talk about the viscosity
solutions of Hamilton--Jacobi, which is joint
interest with the host, ``but we also have
quantum algorithms for other types of
solutions: you can come and ask me, or just
email me''.

\medskip\noindent
{\bf Spoken: viscosity and Hamilton--Jacobi.}
A workable definition of the viscosity solution
is to add an artificial viscosity term, which
does not exist in reality, like a
regularisation in optimisation. Modern schemes
use a nonlinear viscosity, active only where
there are corners, which is efficient. For
quantum one wants linearity, so she goes back
to the artificial viscosity of the 1950s, a
linear diffusion term: as $\nu\to0$ the
solution of the viscous equation tends to the
viscosity solution. Adding a term makes it
possible to linearise. She says that Cole--Hopf
gives a heat equation with a source term, which
is of a form already handled, and that one
cares about $S$, not $u$.

\medskip\noindent
{\bf Spoken: Hamilton--Jacobi and the
Lagrangian.}
Where does the equation come from? In classical
mechanics the action $S$ solves the
Hamilton--Jacobi equation, and this ``hints at
an approach for Lagrangian formulation'':
Newton's equations with the Hamiltonian,
or minimising an action with the Lagrangian,
the two being Legendre--Fenchel dual. This is
the origin of the time-stepping formula
$S^{n+1}=\inf\dots$, whose infimum gives
$-\mathcal{H}$. Why add the entropy regulariser?
Without it the infimum is hard to solve; with
it the minimiser is a Gibbs state (an inverse
exponential), and the minimum value is a free
energy, proportional to minus the logarithm of
a partition function, as in thermodynamics. So
$\nu$ plays the role of a temperature in the
Gibbs weight, and $\gamma$ is a probability
density on the velocities $v$. Cole--Hopf does
not work on the whole equation for a general
Hamiltonian, but it works in each small step,
so one has either a time-stepping algorithm or
a linear PDE. For the observables she says: the
gradient is $\nabla S=\nabla u/u$ up to the
factor $-2\nu$; replacing $d/dx$ by $i\hat p$,
it is obtained by measuring the momentum; the
cost she quotes is of order $d^2/\epsilon$
(the transcript is unclear whether it is
$\epsilon$ or $\epsilon^2$) against
$d^4/\epsilon^2$ for the value of $S$, ``not
optimal'' and improvable.

\medskip\noindent
{\bf Spoken: Wasserstein proximal operators.}
She says the host asked her to talk about them
as an application of the viscosity solution.
Question: which new distribution best reduces a
loss while staying close to the current one?
Different distances give different algorithms.
Gradient descent uses the $L^2$ distance, which
does not preserve normalisation and positivity
of a probability, while the Wasserstein distance
does. In its dynamical form one minimises a
kinetic energy subject to the continuity
equation, which is what makes $\rho$ move as a
probability, with $u$ read as a velocity;
the constraint is a PDE itself with fixed
initial and final distributions. She then adds
a regularising diffusion term, ``a viscosity
term similar to the one in our
Hamilton--Jacobi equation'', so that the
continuity equation becomes the Fokker--Planck
equation, and adds the Lagrange multipliers,
which gives the backward-forward system on the
slides.

\medskip\noindent
{\bf Spoken: the end of the talk.}
(From the transcript of the last three
recordings, which is rough; words in
brackets are my reading of garbled speech.)
Liu summarises the Wasserstein part. To
solve the Hamilton--Jacobi equation one first
adds the viscous term, since the genuine
Hamilton--Jacobi equation is not known to be
solvable quantumly, and then uses the
Cole--Hopf transformation to make the problem
exactly linear. In the forward-backward system
one uses Cole--Hopf again: the final
distribution $\rho(T,x)$ is the
Cole--Hopf factor times $\psi(T,x)$, which comes
from an initial condition, and ``then
everything now is linear''. The two heat
equations are solved one after the other; the
first is solved backwards; there is a
quantum protocol, which the transcript renders
as ``quantum habitar division'' and which
is probably a Hadamard (entrywise) division,
that forms the product of the initial state
with the inverse of the final one. Then the
second heat equation prepares
$\psi(T)$, and from it $\rho(T)$, which encodes
the proximal operator.

\medskip\noindent
{\bf Spoken: the company.}
Liu co-founded a company, which the transcript
names Unitary Lab. Its web page, with an
English version, has a prompt where one
describes the PDE, or one can choose a PDE
from a list, then choose boundary conditions
or initial data on the right. She says there
is no quantum computer yet behind it: the
simulation is classical, with the quantum
circuits in the background. They would
include other people's algorithms if asked.

\medskip\noindent
{\bf Spoken: questions after the talk.}
Three questions, in my summary of the
transcript.
\begin{enumerate}
\item The first questioner asks whether
Hamilton--Jacobi equations solved quantumly
could be connected with search
algorithms, as a strategy for finding where
trajectories or waves can arrive. Liu first
asks what search means here, and ends with: for
a mapping onto quantum algorithms one wants an
algorithm as simple and as linear as possible.
Many search procedures need another protocol
to check that one went to the right place, and
measurement disturbs a quantum state. So
that kind of approach would get messy quickly
and the quantum advantage would be lost. She
adds the point on which the whole talk
insists: if one wants the solution at every
grid point, one loses, because each measurement
destroys the state, so the number of runs is
at least the size of the vector. That is why she
cares about observables: the value at one
point, the minimum value, and so on, which
need one type of measurement.

\item The second questioner asks whether one
can go beyond the linear case, that is, other
Hamiltonians. Liu: very hard. The linear case
works because the two equations of the
forward-backward system decouple: one starts
from a terminal condition, solves backwards,
and puts it into the initial condition. For a
nonlinear problem they become coupled
nonlinearly and she does not know a linear
mapping. One can go beyond the quadratic
Hamiltonian by the time-step method of the
slides, but then one deals with the two
equations one after the other, not together;
and it is only the entropic regularisation
that makes Cole--Hopf linear. She would be
interested in a different regulariser, and
offers to talk about it.

\item The third questioner asks where the
quantum advantage is. Liu: in the simulation
itself. If the PDE is high-dimensional,
finite-difference or similar methods are
exponential in the dimension $D$, and the
quantum algorithm saves that, provided that one
asks for observables whose output does not
cost too much. Here the algorithm only
prepares the quantum state of the proximal
operator; it is a subroutine, from which one
can compute expectation values, or one can
keep taking proximal steps as part of a
Wasserstein flow. In the linear case the
proximal step is exact, up to the
regularisation, and the dependence on the
regularisation parameter is the same as for
the viscosity solution.
\end{enumerate}
The session chair thanked the speaker and the
audience applauded.

\medskip\noindent
{\bf Coverage.}
The photographs cover about 50 of the more
than 100 slides (the deck has at least 118
pages) and no title slide, outline or final
summary slide. The spoken remarks come from
all four recordings, which together cover the
talk and the questions, but the transcripts are
rough: they are my reading of automatic speech
recognition, with garbled names and
formulas, and some stretches (for instance the
numbers of the experiments) could not be
verified. All the mathematics above is as
shown on the slides, plus the editorial
answers, which are mine.

\section{Robust and Efficient Mean
Estimation via Soft Shrinkage}
\label{section-robust-efficient-mean}
%% speaker: Paulo Orenstein, IMPA
%% place: Conference for the 60th anniversary
%% place: of the Department of Mathematics,
%% place: PUC-Rio
%% date: August 12, 2026

\medskip\noindent
{\bf Conference information.}

\href{%
https://sites.google.com/view/%
60thanniversary-dmat-puc-rio/Home
}{International Conference Celebrating
the 60th Anniversary of the Department
of Mathematics at PUC-Rio}

August 12--14, 2026,
RDC Auditorium, PUC-Rio,
Rio de Janeiro.

\medskip\noindent
{\bf Speaker.}

Paulo Orenstein, IMPA.

\medskip\noindent
{\bf Date and time.}

August 12, 2026, 11:15--12:15.

\medskip\noindent
{\bf Title.}

Robust and Efficient Mean Estimation
via Soft Shrinkage.

\medskip\noindent
{\bf Abstract.}

How should we weight observations when
estimating a mean? The sample mean is
efficient and minimax but fragile under
heavy tails; robust estimators resist
heavy tails but typically sacrifice
efficiency. We show a soft-shrinkage
calibration improves on both fronts:
observations are smoothly downweighted
by their distance from a pilot estimate,
with the shrinkage scale calibrated so
that the total rejected weight equals a
prescribed budget. Any reasonable pilot
is upgraded to sub-Gaussian concentration
with the minimax constant $\sqrt{2}$ under
finite variance alone, and the budget
doubles as an exact breakdown point. The
mechanism trades variance for bias, and
the bias can be quantified exactly. This
yields a second-order theory establishing
an MSE expansion that identifies when
shrinkage beats the sample mean,
Berry--Esseen and moderate deviation
approximations, and asymptotically exact
confidence intervals. Beyond the theory,
the results give concrete guidance for
designing mean estimators and bridge
modern sub-Gaussian estimation with
classical robust statistics. This is
joint work with Ant\^onio Cat\~ao.

\medskip\noindent
{\bf Notes.}
The talk starts introducing the mean
$\mu$ of a sample $X_i$, which is 
the sum of all the variables divided
by the size of the sample. This is
an efficient estimate, but not robust.
Robustness can be explained with atypical
data: suppose that you have a very atypical
datum in the sample. This will move the
mean $\mu$ quite a lot, but this isn't very
good since it was just one extraordinary
datum. Then we introduce the trimmed mean,
which uses a new variable to trim, or
remove, the tails of the data, i.e.
the most extraordinary data points
to avoid the case we just explained.
The trimmed mean is robust, but looses
efficiency since we need to chop off a 
piece of our data, which isn't ideal.

The talk aims at introducing a new 
estimate which lies in between the 
usual mean and the trimmed mean.
The idea is to introduce weights to every
data using a variable  $\alpha$ (which
depends on the whole data).
Such a new estimate is both robust
and efficient.

Developing these kinds of statistical
gadgets has applications in mathematics,
science, computer science, artificial
intelligence, etc. In particular, this can
be thought of as a more sophisticated method
than simple gradient descent used in machine
learning methods, the one used in the AI
course at EGAXII.

\section{Recovering Structure from Binary
Graphical Model Dynamics}
\label{section-recovering-binary-graphical}
%% speaker: Roberto Imbuzeiro, IMPA
%% place: Conference for the 60th anniversary
%% place: of the Department of Mathematics,
%% place: PUC-Rio
%% date: August 12, 2026

\medskip\noindent
{\bf Conference information.}

\href{%
https://sites.google.com/view/%
60thanniversary-dmat-puc-rio/Home
}{International Conference Celebrating
the 60th Anniversary of the Department
of Mathematics at PUC-Rio}

August 12--14, 2026,
RDC Auditorium, PUC-Rio,
Rio de Janeiro.

\medskip\noindent
{\bf Speaker.}

Roberto Imbuzeiro, IMPA.

\medskip\noindent
{\bf Date and time.}

August 12, 2026, 16:00--17:00.

\medskip\noindent
{\bf Title.}

Recovering structure from binary
graphical model dynamics.

\medskip\noindent
{\bf Abstract.}

We report on work in progress with
Guilherme Ost (UFRJ) and Guilherme Reis
(UFF) on a class of toy models for
neuronal dynamics. These models consist
of sets of units taking $0/1$ values,
which interact pairwise via a matrix of
weights. The main question we address is
how much of the weight structure can be
recovered from the dynamics. We will
focus on ``mean field'' cases where
individual pairwise interactions may be
quite weak. Interestingly, this property
turns out to be helpful.

\medskip\noindent
{\bf Notes.}
This talk made me remember Carina Curto's
neural networks toy models. It looks
like they worked in (oversimplified)
models describing neural networks.

(AI paragraph)
The two approaches are complementary.
Curto starts with a directed connectivity
graph and studies the resulting dynamics,
such as fixed points and attractors.
Imbuzeiro starts with observed binary
dynamics and asks what can be recovered
about the unknown interaction graph. Thus
Curto studies structure to dynamics, while
Imbuzeiro studies dynamics to structure.


A paper mentioned in the talk is
\href{%
https://doi.org/10.1214/25-AOS2592
}{Inferring the dependence graph density
of binary graphical models in high
dimension}, by Julien Chevallier,
Eva L\"ocherbach, and Guilherme Ost. It
appeared in \emph{The Annals of
Statistics}, volume 54, number 2 (2026),
pages 861--881.

\section{Detecção de Comunidades para
Modelos Gráficos Binários de Alta Dimensão}
\label{section-ost-colmea-2026}
%% speaker: Guilherme Ost, IMPA
%% place: COLMEA, IMPA
%% date: September 23, 2026

\noindent
Guilherme Ost (IMPA).
September 23, 2026.
COLMEA, IMPA, 15:40--17:00.

\medskip\noindent
{\bf Abstract.}
Consideramos um sistema de $N$ unidades
binárias que interagem aos pares por meio
de uma matriz de pesos. Esta matriz é
possivelmente aleatória e possui uma
estrutura em blocos, codificando conjuntos
de unidades que desempenham um papel
semelhante no sistema, chamados de
comunidades. Tanto a matriz de pesos quanto
essas comunidades são desconhecidas. Nesta
palestra, focarei na seguinte questão
estatística: ao observar apenas os estados
das $N$ unidades ao longo de $T$ passos de
tempo, conseguimos recuperar as diferentes
comunidades do sistema? Nosso objetivo é
apresentar garantias teóricas de alguns
métodos estatísticos simples que nos
permitem recuperar essas comunidades a
partir dos dados observados, desenvolvidos
recentemente com J. Chevallier (Grenoble,
França).
\medskip\noindent
{\bf Notes.}
The classical version of this problem is
``community detection in networks.'' The
simplest case is described by a symmetric
stochastic block model.

Let the $N$ vertices be divided into two
communities
$
\mathcal{C}_1,\mathcal{C}_2
$
of equal size. Every pair $(i,j)$ of
vertices is connected independently with
probability $p$ if $i\sim j$, meaning that
$i$ and $j$ belong to the same community,
and with probability $q$ if $i\not\sim j$.
We suppose
$
p>q>0.
$

The resulting connectivity is represented
by an adjacency matrix. The problem is:
given a permuted version of the adjacency
matrix, say $\theta$, for an unknown
permutation of the vertices, recover the
communities $\mathcal{C}_1$ and
$\mathcal{C}_2$.

There is a solution to this classical
problem using spectral methods.

Here $p$ is the probability of connecting
two vertices in the same community, while
$q$ is the probability of connecting two
vertices in different communities.

The basic quantitative question is how
separated $p$ and $q$ must be in order to
recover the communities. Equivalently, how
small can $p-q$ be before the community
structure becomes statistically invisible?

\medskip\noindent
{\bf Some classical thresholds.}
Yun and Prouti\`ere (2014) use the sparse
parametrization
$
p=\frac{a_N}{N},
\qquad
q=\frac{b_N}{N}.
$
Take a vertex $v\in\mathcal{C}_1$.
Since $|\mathcal{C}_1|=N/2$, the expected
number of its neighbours inside
$\mathcal{C}_1$ is approximately
$
|\mathcal{C}_1|p
=\frac{N}{2}\frac{a_N}{N}
=\frac{a_N}{2}.
$
Likewise, the expected number of
neighbours of $v$ in $\mathcal{C}_2$ is
$
|\mathcal{C}_2|q=\frac{b_N}{2}.
$
Thus $a_N-b_N$ measures the separation
between the two expected degree profiles.

Abbe, Bandeira and Hall (2016) consider
$
p=\frac{a\log N}{N},
\qquad
q=\frac{b\log N}{N}.
$
They obtain a sharp threshold for exact
recovery. Above
$
\frac{a+b}{2}-\sqrt{ab}>1
$
exact recovery is possible with
probability tending to one; below the
corresponding strict inequality it is
impossible. Equivalently, the threshold is
$
(\sqrt a-\sqrt b)^2=2.
$

\medskip\noindent
{\bf From observing the graph to observing
binary data.}
The problem of interest in this talk is
harder: the graph itself is not observed.
The motivating example comes from
neuroscience and spike-train data. A
``spike'' is an action potential, i.e. an
event in which a neuron fires.

Discretize time into windows and record
$1$ if neuron $i$ fires during a window
and $0$ otherwise. At each time $t$ this
gives a binary state
$
X_t\in\{0,1\}^N.
$
Over $T$ time windows one therefore
observes a binary data matrix. The hidden
question is whether one can reconstruct
the excitatory and inhibitory communities
from these data alone. This is the
community detection problem from binary
data.

\medskip\noindent
{\bf Ingredients.}
\begin{itemize}
\item
An unobserved graph $\theta$, possibly
random, directed, and with block
structure.

\item
A probability law $\mathbb{P}_\theta$ on
trajectories of binary network states.

\item
An observed sample
$
(X_1,\ldots,X_T)\sim\mathbb{P}_\theta,
$
from which we want to recover the
communities.
\end{itemize}

\medskip\noindent
{\bf IA-remark.}
Morally, a trajectory is just the movie of
the network through time. Each frame
$X_t\in\{0,1\}^N$ records which of the
$N$ neurons fired in the $t$th time
window. Thus
$
(X_1,\ldots,X_T)
$
is one complete observed history of the
network. A probability law on trajectories
says how likely the different possible
movies are.

\medskip\noindent
{\bf Independent benchmark.}
In an i.i.d. model one would have
$
\mathbb{P}_\theta
=\bigotimes_{t=1}^T P_\theta,
$
where $P_\theta$ is a probability measure
on $\{0,1\}^N$. A famous example is the
Ising blockmodel of Berthet, Rigollet and
Srivastava (2019), where independent
binary observations are used to recover
an underlying block structure.

\medskip\noindent
{\bf The speaker's model.}
The model goes beyond the i.i.d.
assumption. The vectors
$
X_1,\ldots,X_T\in\{0,1\}^N
$
are dependent under
$\mathbb{P}_\theta$. The random graph has
directed interactions and a block
structure. A connection emanating from an
excitatory neuron has positive weight.

Under $\mathbb{P}_\theta$, the process
$
(X_t)_{1\leq t\leq T}
$
is a stationary Markov chain.

The statistical problem is: from the
sample $X_1,\ldots,X_T$, recover the
partition
$
(\mathcal{P}_+,\mathcal{P}_-).
$
The model parameters $\mu$, $\lambda$ and
$\theta$ are unknown. The community sizes
have the form
$
|\mathcal{P}_+|=Nr_+,
\qquad
|\mathcal{P}_-|=Nr_-,
$
where
$
0<r_+,r_-<1,
\qquad
r_++r_-=1,
$
and $r_+$ and $r_-$ are also unknown.

\medskip\noindent
{\bf AI-reference check.}
Chevallier and Ost's paper
``Community detection for binary graphical
models in high dimension'' studies this
dependent-data problem. Their exact
recovery probability tends to one in the
regime
$
\frac{N}{\sqrt{T}}\log(NT)\longrightarrow0.
$

\medskip\noindent
{\bf AI-speaker profile.}
Guilherme Ost is a professor at UFRJ who
works in probability and stochastic
processes. His research repeatedly studies
how hidden structure can be inferred from
dependent stochastic observations. Topics
include neuronal interaction graphs from
spike trains, high-dimensional binary
graphical models, and random walks on
random graphs. In 2026 he is also teaching
High-Dimensional Statistics at IMPA.

\bibliography{my}
\bibliographystyle{amsalpha}

\end{document}
